% Dobrushin uniqueness for the four-dimensional SU(2) Wilson lattice gauge measure for beta < 2/9.
% Version 1.4, 28 October 2026 (literature update: [BB26]; no change to the mathematics of v1.3). Compile: pdflatex (three passes). Cyrillic: T2A, family lmr substituted by cmr
% (cm-super Type 1 fonts), for the verbatim quotations of Dobrushin (1970).
\documentclass[11pt,a4paper]{article}
\usepackage[T2A,T1]{fontenc}
\usepackage[utf8]{inputenc}
\usepackage[russian,english]{babel}
% Unicode characters occurring in the verbatim Lean statements (Sections 4.4 and 5.2)
\DeclareUnicodeCharacter{211A}{\ensuremath{\mathbb{Q}}}
\DeclareUnicodeCharacter{211D}{\ensuremath{\mathbb{R}}}
\DeclareUnicodeCharacter{2264}{\ensuremath{\le}}
\DeclareUnicodeCharacter{2208}{\ensuremath{\in}}
\DeclareUnicodeCharacter{03B2}{\ensuremath{\beta}}
\DeclareUnicodeCharacter{03BA}{\ensuremath{\kappa}}
\DeclareUnicodeCharacter{2200}{\ensuremath{\forall}}
\DeclareUnicodeCharacter{2192}{\ensuremath{\to}}
\DeclareUnicodeCharacter{2082}{\ensuremath{_2}}
\usepackage{lmodern}
\DeclareFontFamilySubstitution{T2A}{lmr}{cmr}% Latin Modern has no T2A shapes; use cm-super for Cyrillic
\usepackage{amsmath,amssymb,amsthm,mathrsfs}
\usepackage{booktabs,tabularx}
\usepackage{geometry}
\geometry{margin=2.6cm}
\usepackage[colorlinks=true,linkcolor=blue,citecolor=blue,urlcolor=blue]{hyperref}
\usepackage{xurl}% \path{} for hashes and file paths, breakable anywhere

\hypersetup{pdftitle={Dobrushin uniqueness for the four-dimensional SU(2) Wilson lattice gauge measure for beta < 2/9, with interval/Lean certificates},pdfauthor={K\'evin R\'emondi\`ere}}
\newtheorem{theorem}{Theorem}
\newtheorem{proposition}[theorem]{Proposition}
\newtheorem{lemma}[theorem]{Lemma}
\theoremstyle{remark}
\newtheorem{remark}[theorem]{Remark}
\theoremstyle{definition}
\newtheorem{definition}[theorem]{Definition}

\newcommand{\R}{\mathbb{R}}
\newcommand{\Z}{\mathbb{Z}}
\newcommand{\E}{\mathbb{E}}
\newcommand{\Cov}{\operatorname{Cov}}
\newcommand{\Var}{\operatorname{Var}}
\newcommand{\ReTr}{\operatorname{Re}\operatorname{Tr}}
\newcommand{\st}[1]{\textsf{[#1]}}
\newcommand{\ru}[1]{\foreignlanguage{russian}{#1}}
\newcommand{\Zb}{\bar Z}
\newcommand{\Dch}{D_{\mathrm{ch}}}
\newcommand{\rch}{\rho_{\mathrm{ch}}}

\title{Dobrushin uniqueness for the four-dimensional\\ SU(2) Wilson lattice gauge measure for $\beta<2/9$,\\
with interval/Lean certificates\\[4pt]
\large Version 1.4}
\author{K\'evin R\'emondi\`ere\thanks{Independent researcher, Oloron-Sainte-Marie, France. ORCID
\href{https://orcid.org/0009-0008-2443-7166}{0009-0008-2443-7166}. The note, the certificate code and the Lean
files were prepared with the assistance of Claude (Opus 5.5, Anthropic); every claim carries one of the status
labels defined below.}}
\date{28 October 2026 (v1.4: literature update, \S\ref{sec:lit}; v1.3 of 28 September 2026, not publicly released)}

\begin{document}
\maketitle

\begin{abstract}
We write out Dobrushin's uniqueness criterion, in its Wasserstein ($W_1$) form, for the $SU(2)$ Wilson lattice
gauge measure on $\Z^4$, and compute the single-link influence coefficient \emph{exactly}, uniformly in the
field, for two metrics on $SU(2)\cong S^3$.
For the Euclidean (chordal) metric of $\R^4$, the metric used by Levin (1980), we show that the one-site
coefficient in a field of strength $\kappa\le4/3$ equals $\E_\kappa[x_\perp^2]$, the second moment of a coordinate
orthogonal to the field; it is maximal at zero field, where it equals $1/4$. Hence each of the $18$ neighbouring
links has influence exactly $\beta/4$, the Dobrushin sum is $9\beta/2$, and the (DLR) Gibbs measure exists and is
unique for $\beta<2/9\approx0.2222$, in the normalisation $\beta\sum_p(1-\tfrac12\ReTr U_p)$.
This confirms, for $N=4$ and $|g|\le4/3$, the maximum-at-zero-field property conjectured by Levin (Remark~3.7), which
Bhattacharjee and Bhattacharya prove for $N=4$ and all fields \cite[Theorem~7]{BB26};
accordingly, $2/9$ is exactly Levin's published bound transferred to the Wilson link, $2\sqrt5/45\approx0.0994$,
with his factor $\sqrt5$ replaced by $1$.
The proof bounds the covariance by the $L^1$ norm of an explicit divergence-free vector field on the ball with
prescribed normal flux, eliminates the direction analytically, and reduces the claim to a one-variable
inequality on $[0,4/3]$, proved by interval arithmetic (two independent implementations); its final rational
step is checked in Lean~4.
For the geodesic metric (the main result of v1.2 of this note) the exact coefficient is $8\beta/(9\pi)$,
uniformly for $\kappa\le3\pi/8$, which gives the weaker range $\beta<\pi/16\approx0.196$ by the same method. Each
threshold is sharp for its criterion.
The single-site machinery is not new: Levin (1980, Euclidean metric) and Royer (1986, geodesic metric) treat
rotator models, and Levin's many-body version already covers the Wilson model; none of the cited papers
mentions gauge theories. For comparison, the explicit range of Shen--Zhu--Zhu is $\beta<1/12$, the
total-variation criterion (Gross 1979, Simon 1979) gives $\beta<1/18$ transferred, and global Bakry--\'Emery on
link variables cannot pass $\beta=1/8$. No statement on decay of correlations, spectral gap or log-Sobolev
inequality is claimed.
\emph{Prior public statement.} The uniqueness range $\beta<2/9$ for the four-dimensional $SU(2)$ model, with
exponential decay and a lattice mass gap, and the all-field bound $W_1(\nu_b,\nu_{b'})\le\frac14|b-b'|$ on $S^3$, were
stated publicly, with proofs, in the preprint \cite{BB26} of 18 October 2026 (Theorems~E, F and~7 there). This note
was prepared independently and makes \emph{no priority claim}. Its own contributions are the exact identification
of the one-link chordal coefficient as $\E_\kappa[x_\perp^2]$ for $\kappa\le4/3$, with the explicit remainder
$\frac14-\E_\kappa[x_\perp^2]\ge\kappa^2/(96\,\Zb(\kappa))$, the exact geodesic coefficient $8\beta/(9\pi)$ for
$\kappa\le3\pi/8$, two independent interval-arithmetic certificates, Lean-checked final steps, and an explicit
transfer of the published constants of Levin, Royer and Shen--Zhu--Zhu with their conventions.
\end{abstract}

\paragraph{Status labels.}
Every claim carries one label:
\st{proved} (complete argument written here or in the cited source);
\st{interval-certified} (rigorous modulo trust in the interval library and in the reductions of
\S\S\ref{sec:chfield}--\ref{sec:chalpha} and \S\ref{sec:reduction}, which are proved but not formalised);
\st{Lean-checked} (kernel-checked in Lean~4, axioms \texttt{propext}, \texttt{Classical.choice}, \texttt{Quot.sound} only;
this label never covers more than the Lean statements quoted verbatim in \S\ref{sec:chcap} and \S\ref{sec:cap},
which are rational inequalities, conditional implications and threshold arithmetic);
\st{read} (statement read in the original source, location given);
\st{conjecture}.
The label \st{numerical} marks floating-point evidence which is \emph{not} part of any proof.

\section{Statement and conventions}\label{sec:statement}

Let $E$ be the set of (positively oriented) links of $\Z^4$ and $\mathcal P$ the set of unoriented plaquettes.
A configuration is $U=(U_\ell)_{\ell\in E}\in SU(2)^E$; $U_p$ is the ordered product around $p$ (its
real trace does not depend on the starting point or orientation, since $\ReTr W=\ReTr W^{-1}$ on $SU(2)$).
For a finite $\Lambda\subset E$ and boundary condition $\bar U$, the Wilson specification is
\begin{equation}\label{eq:wilson}
  \mu^{\bar U}_{\Lambda}(dU_\Lambda)\;\propto\;
  \exp\Big(-\beta\sum_{p\cap\Lambda\neq\emptyset}\big(1-\tfrac12\ReTr U_p\big)\Big)\prod_{\ell\in\Lambda}dU_\ell ,
\end{equation}
with $dU$ the normalised Haar measure. Throughout, $\beta=\beta_W$ is the \emph{Wilson} coupling, and
the bare coupling $g$ is \emph{defined} by $\beta_W=2N/g^2$, i.e.\ $\beta_W=4/g^2$ for $SU(2)$; the coupling
$g_0$ of Osterwalder--Seiler is a different normalisation (Table~\ref{tab:conv}). A Gibbs measure is a probability measure on $SU(2)^E$ satisfying the
DLR equations for \eqref{eq:wilson}.
We identify $SU(2)$ with the unit sphere $S^3\subset\R^4$ (unit quaternions) and use two metrics on it: the
\emph{chordal} (Euclidean) distance $\rch(x,y)=|x-y|$ of $\R^4$, of diameter $2$, and the \emph{geodesic}
(great-circle) distance $\rho$, of diameter $\pi$. They generate the same topology and Borel $\sigma$-algebra.

\begin{table}[htbp]
\centering\small
\begin{tabular}{lll}
\toprule
convention & weight & relation to $\beta_W$ (for $SU(2)$)\\
\midrule
Wilson (this note) & $\exp(-\beta_W\sum_p(1-\frac1N\ReTr U_p))$ & ---\\
$\frac{\beta}{N}\ReTr$ & $\exp(\frac{\beta}{N}\sum_p \ReTr U_p)$ & $\beta_W=\beta$\\
raw trace & $\exp(\beta\sum_p\ReTr U_p)$ & $\beta_W=2\beta$\\
't~Hooft (SZZ, CNS) & $\exp(N\beta\sum_p\ReTr U_p)$ & $\beta_W=N^2\beta=4\beta$\\
Osterwalder--Seiler & $\exp(\frac{1}{2g_0^2}\sum_P \operatorname{Re}\chi(g_P))$ & $\beta_W=1/g_0^2$ (hypothesis${}^\ast$), i.e.\ $g_0=g/2$\\
\bottomrule
\end{tabular}
\caption{Conventions. The 't Hooft line is \cite[(1.2)]{SZZ23} (see \S\ref{sec:lit}); the
Osterwalder--Seiler line is \cite[(2.1)--(2.2), p.~442]{OS78}, read on the page image
(\st{read}). ${}^\ast$The conversion $\beta_W=1/g_0^2$ is a \emph{hypothesis}: it holds if $\chi$ is the
fundamental character and each geometric plaquette is counted once in $\sum_P$. On p.~442, $\chi$ is
``a character on $\mathcal G$'' and plaquettes are ``closed loops $uvxy$''; neither condition is stated
there. The conversion is used only to translate thresholds, never in a proof.}
\label{tab:conv}
\end{table}

\begin{theorem}[Main result]\label{thm:main}
For $0\le\beta<2/9\approx0.2222$ there is exactly one Gibbs measure for the $SU(2)$ Wilson
specification \eqref{eq:wilson} on $\Z^4$.
Equivalently: $g^2>18$ with $g$ defined by $\beta_W=2N/g^2=4/g^2$; in the 't~Hooft normalisation of
\cite{SZZ23}, $\beta_{\mathrm{tH}}<1/18\approx0.0556$; in the normalisation of \cite{OS78} (under the hypothesis
of Table~\ref{tab:conv}), $g_0^2>9/2$.
\end{theorem}
\noindent Status: existence \st{proved} (Proposition~\ref{prop:exist});
uniqueness \st{read} (Dobrushin 1970, Thm~4) + \st{proved} (\S\ref{sec:influence}, chordal metric) +
Theorem~\ref{thm:chordal}, which is \st{proved} + \st{interval-certified}; its last rational step and the
threshold arithmetic are \st{Lean-checked} (\S\ref{sec:chcap}).

\begin{theorem}[Geodesic criterion; the main result of v1.2]\label{thm:geo}
With the geodesic metric $\rho$, Dobrushin's sum for \eqref{eq:wilson} equals $16\beta/\pi$ for
$0\le\beta<\pi/16\approx0.19635$; so the geodesic $W_1$ criterion gives uniqueness for $\beta<\pi/16$ and fails at
$\beta=\pi/16$.
\end{theorem}
\noindent Status: \st{proved} (\S\ref{sec:influence}) + Theorem~\ref{thm:onelink}, which is \st{proved} for
$\kappa=0$ and \st{interval-certified} for $\kappa>0$, with its last rational step \st{Lean-checked}
(\S\ref{sec:cap}). The uniqueness statement for $\beta<\pi/16$ is contained in Theorem~\ref{thm:main}; what
Theorem~\ref{thm:geo} adds is the exact geodesic coefficient $8\beta/(9\pi)$ and its field-uniformity.

Each threshold is \emph{sharp for its criterion}: with single-site conditionals, the chordal sum equals $9\beta/2$
and the geodesic sum $16\beta/\pi$ exactly (Proposition~\ref{prop:influence}), so the chordal criterion fails at
$\beta=2/9$ and the geodesic one at $\beta=\pi/16$. Nothing is claimed about other metrics, block or multi-site
versions of the criterion, or the true uniqueness region.

\section{Dobrushin's theorem and its hypotheses}\label{sec:dobrushin}

We use the original Russian article \cite{Dob70} (Math-Net.Ru scan \texttt{tvp1856}; the
quotations below were read on the rendered page images; English wording is ours).

\paragraph{Setting (\S1, p.~469).} $T$ is a countable set; the random variables $\xi_t$ take values in a
measurable space $(X,\mathfrak B)$. Dobrushin's standing assumption is that all probability measures on
finite products of $(X,\mathfrak B)$ are perfect (``\ru{все вероятностные меры на прямых произведениях
конечного числа экземпляров измеримого пространства $(X,\mathfrak B)$ совершенны}''), and for this it
suffices (verbatim) that
``\ru{$\mathfrak B$ было $\sigma$-алгеброй всех борелевских подмножеств для некоторой метрики,
превращающей $X$ в полное сепарабельное метрическое пространство}''.
Given an integer function $\nu$ on $T$,
\begin{equation*}
G_{t_0}=\{t:\nu(t)\le\nu(t_0),\ t\neq t_0,\ t\in T\},\qquad t_0\in T \qquad\text{\cite[(1.2)]{Dob70}},
\end{equation*}
and a field has conditional distributions $\{p_{t_0}\}$ if
$\mathbf P\{\xi_{t_0}\in B\mid\xi_t,t\in G_{t_0}\}=p_{t_0}(B\mid x_t,t\in G_{t_0})$ \cite[(1.3)]{Dob70}.

\paragraph{Wasserstein distance (\S3, p.~479).}
``\ru{Назовем \emph{расстоянием Васерштейна} (см.~[12]) число}
$R(P,Q)=\inf \mathrm{M}\rho(\eta,\zeta)$, (3.1)
\ru{где нижняя грань берется по всем парам случайных величин $\eta,\zeta$ \dots\ таким, что $\eta$ имеет
распределение $P$ и $\zeta$ --- распределение $Q$}''. This is $W_1$ defined by couplings, for the metric
$\rho$ of $X$ (any metric making $X$ complete and separable).

\paragraph{Theorem 4 (\S5, p.~488), verbatim.}
``\ru{Для того чтобы существовало не более, чем одно, случайное поле $\{\xi_t, t\in T\}$ с заданными
условными распределениями $\{p_{t_0}(\cdot|\cdot), t_0\in T\}$, для которого при некоторых $x_0\in X$ и
$C<\infty$}
$\mathrm{M}\rho(\xi_t,x_0)\le C<\infty$, $t\in T$, (5.1)
\ru{достаточно, чтобы при некотором $d<1$ и всех $t_0\in T$}
$\sum_{t\in G_{t_0}} r(t,t_0)\le d<1$, (5.2)
\ru{где числа $r(t,t_0)$ \dots\ таковы, что расстояние Васерштейна}
$R(p_{t_0}(\cdot|x^1_t,t\in G_{t_0}),p_{t_0}(\cdot|x^2_t,t\in G_{t_0}))\le\sum_{t\in G_{t_0}}r(t,t_0)\rho(x^1_t,x^2_t)$ (5.3)
\ru{\dots\ Условие (5.3) выполнено, в частности, если}
$r(\bar t,t_0)=\sup R(\cdots)/\rho(x^1_{\bar t},x^2_{\bar t})$, (5.4)
\ru{где верхняя грань берется по всем парам условий \dots, отличным лишь при $t=\bar t$}.''
(Translation: at most one field with the given conditionals and property (5.1) exists provided that, for
some $d<1$ and all $t_0$, (5.2) holds with numbers $r$ satisfying the Wasserstein bound (5.3); (5.3)
holds in particular for the supremum (5.4) over pairs of conditions differing only at $\bar t$.)

\paragraph{Theorem 1 (existence, pp.~472--473).} Two conditions suffice; on p.~473 Dobrushin writes
``\ru{Условие 1 заведомо выполнено, если $X$ --- бикомпакт. (Можно положить $h(x)\equiv1$.)}'' and, in the
Markov case (conditionals depending on finitely many $x_t$, $t\in U_{t_0}$), ``\ru{условие 2 означает
непрерывность $p_{t_0}(\cdot\mid x_t, t\in U_{t_0})$ по переменным $x_t$ в топологии слабой
сходимости}''.

\paragraph{Verification of the hypotheses for \eqref{eq:wilson}.} Let $d$ be either $\rch$ or $\rho$.
\begin{enumerate}
\item $T=E$ is countable; $X=SU(2)\cong S^3$ with $d$ is compact, hence complete and separable;
$\mathfrak B$ is its Borel $\sigma$-algebra (the same for both metrics). \st{proved}
\item Take $\nu\equiv0$; then $G_{t_0}=T\setminus\{t_0\}$ (the case Dobrushin himself mentions on
p.~488, ``\ru{В случае $G_{t_0}=\{t:t\neq t_0\}$}''). The kernel $p_{\ell}(\cdot\mid U_{E\setminus\ell})$ is
the single-link conditional law \eqref{eq:q} below; it has a density continuous in all variables, hence
is jointly measurable. Every DLR measure for \eqref{eq:wilson} satisfies \cite[(1.3)]{Dob70} with these
kernels. So uniqueness in Dobrushin's class implies uniqueness of the Gibbs measure. \st{proved}
\item (5.1) holds for every field with $x_0=\mathbf 1$ and $C=\operatorname{diam}S^3$ ($2$ for $\rch$, $\pi$ for $\rho$). \st{proved}
\item (5.2)--(5.4): $p_\ell$ depends only on the $18$ neighbouring links (\S\ref{sec:influence}), so
$r(t,\ell)=0$ for all other $t$. With $r$ the supremum (5.4): for $d=\rch$ and $\beta<2/9$,
$\sum_t r(t,\ell)=9\beta/2$ (Proposition~\ref{prop:influence} and Theorem~\ref{thm:chordal}); take $d=9\beta/2<1$.
For $d=\rho$ and $\beta<\pi/16$, $\sum_t r(t,\ell)=16\beta/\pi$ (Theorem~\ref{thm:onelink}). \st{proved} given
Theorems~\ref{thm:chordal} and~\ref{thm:onelink}.
\item Existence: Theorem~1 applies (Condition~1 because $X$ is compact; Condition~2 because the
conditional density is continuous in the neighbouring links), but it yields a field with the
prescribed \emph{single-site} conditionals. We do not use it: existence of a measure satisfying the DLR
equations for \emph{every} finite $\Lambda$ is proved directly in Proposition~\ref{prop:exist}. \st{proved}
\end{enumerate}
A second modern statement of the $W_1$ criterion on complete metric single-spin spaces is
\cite[Thm~3.5]{DS25} (p.~8 of arXiv v4) \st{read}.

\begin{proposition}[Existence in the DLR sense]\label{prop:exist}
For every $\beta\ge0$ there is at least one probability measure $\mu$ on $SU(2)^E$ such that
$\mu\gamma_\Delta=\mu$ for every finite $\Delta\subset E$, where $\gamma_\Delta(\cdot\mid\bar U)=
\mu^{\bar U}_\Delta\otimes\delta_{\bar U_{E\setminus\Delta}}$ is the kernel \eqref{eq:wilson}.
\end{proposition}
\noindent Status: \st{proved}.
\begin{proof}
Write $H_\Delta(U)=\sum_{p\cap\Delta\ne\emptyset}(1-\tfrac12\ReTr U_p)$, a finite sum of continuous functions of
finitely many links, with $0\le H_\Delta\le h_\Delta:=2\,|\{p:p\cap\Delta\ne\emptyset\}|$. Thus
\[
\gamma_\Delta f(\bar U)=Z_\Delta(\bar U)^{-1}\int f(U_\Delta,\bar U_{E\setminus\Delta})\,e^{-\beta H_\Delta(U_\Delta,\bar U)}\,dU_\Delta,
\qquad Z_\Delta\ge e^{-\beta h_\Delta}>0 .
\]
\emph{(a) Feller property.} If $f$ is continuous on the compact space $SU(2)^E$ (product topology), the
integrand is jointly continuous in $(U_\Delta,\bar U)$ and bounded, so $\gamma_\Delta f$ and $Z_\Delta$ are
continuous by dominated convergence, and $\gamma_\Delta f$ is continuous.
\emph{(b) Consistency.} For finite $\Delta\subset\Lambda$ one has $\gamma_\Lambda\gamma_\Delta=\gamma_\Lambda$: the density of
$\gamma_\Lambda(\cdot\mid\bar U)$ is $e^{-\beta H_\Lambda}/Z_\Lambda$ and $H_\Lambda=H_\Delta+H'$ with $H'$ independent of
$U_\Delta$ (the plaquettes meeting $\Lambda$ but not $\Delta$), so by Fubini integrating $\gamma_\Delta f$ against
$\gamma_\Lambda$ first in $U_\Delta$ reproduces $\gamma_\Lambda f$.
\emph{(c) Limit.} Fix $\bar U$ and finite $\Lambda_n\uparrow E$, and let $\mu_n=\gamma_{\Lambda_n}(\cdot\mid\bar U)$. By
compactness (Prokhorov, or Riesz--Markov on $C(SU(2)^E)$) a subsequence converges weakly to some $\mu$.
For finite $\Delta$ and continuous $f$, (b) gives $\mu_n(\gamma_\Delta f)=\mu_n(f)$ as soon as
$\Delta\subset\Lambda_n$; by (a) both sides converge along the subsequence, so $\mu(\gamma_\Delta f)=\mu(f)$ for all
continuous $f$, hence $\mu\gamma_\Delta=\mu$.
\end{proof}
Every such $\mu$ satisfies in particular \cite[(1.3)]{Dob70} with $G_{t_0}=T\setminus\{t_0\}$ (take $\Delta=\{\ell\}$),
which is the direction used for uniqueness in item~2.

\section{The single-link conditional law and the influence matrix}\label{sec:influence}

\paragraph{Quaternion form.} For unit quaternions $u,w$, $\tfrac12\ReTr(uw)=\langle u,\bar w\rangle_{\R^4}$.
Link $\ell$ lies in $2(d-1)=6$ plaquettes; for each such $p$, $\tfrac12\ReTr U_p=\tfrac12\ReTr(U_\ell W_p)$
where the \emph{staple} $W_p$ is a product of the three other links of $p$ (or of their inverses; we use
$\ReTr A=\ReTr A^{-1}$ to put $U_\ell$ first). Hence
\begin{equation}\label{eq:q}
  p_\ell(du\mid\cdot)=q_s(du)\propto e^{\beta\langle u,s\rangle}\,du,\qquad
  s=\sum_{p\ni\ell}\overline{W_p}\in\R^4,\quad |s|\le 6 .
\end{equation}
Write $\kappa=\beta|s|\in[0,6\beta]$, $\hat s=s/|s|$, and $q_\kappa\propto e^{\kappa\langle u,\hat s\rangle}$.
In the notation of Royer \cite[p.~31]{Roy86}, $q_s=\nu(\beta s)$ on $S^n$ with $n=3$: ``For each vector
$h\in\R^{n+1}$ let $\nu(h)$ be the probability measure on $S^n$, the density of which is proportional to
$x\to\exp(x.h)$'' \st{read}, on the Numdam page image; in Levin's notation \cite[(3.11), p.~68]{Lev80} it is
$\mu_{\mathbf h}$ with $\mathbf h=\beta s$ and $N=4$ \st{read}. The single-site analysis below is therefore an
instance of the setting of Levin \cite{Lev80} and Royer \cite{Roy86} for $n$-vector (rotator) models; what
is specific to the gauge model is the staple structure and the count of neighbours. In fact the whole
Wilson model is a $4$-vector model with four-body interactions in the sense of Levin
\cite[Remark~3.6, (3.48), p.~73]{Lev80}: writing $u_1,\dots,u_4\in S^3\subset\R^4$ for the links of $p$ (with
$u^{-1}=\bar u$), $\tfrac12\ReTr U_p=\langle u_1u_2\bar u_3\bar u_4,\mathbf 1\rangle$ is a multilinear function of
$(u_1,u_2,u_3,u_4)$. \st{proved}

\paragraph{Why 18 neighbours.} Each of the $6$ plaquettes through $\ell$ contains $3$ further links.
Two distinct plaquettes through $\ell$ span different coordinate planes, or the same plane on opposite
sides of $\ell$, so they share no link other than $\ell$. Hence the $6\times3=18$ links are pairwise
distinct, each shares exactly one plaquette with $\ell$, and a change of one of them changes exactly one
staple. In general dimension the count is $3\cdot2(d-1)$. \st{proved}

\begin{definition}\label{def:D}
For $\kappa\ge0$ and a unit vector $v\in\R^4$ put
\begin{align*}
\Dch(\kappa,v)&=\sup\{\Cov_{q_\kappa}(f,\langle u,v\rangle): f\ \text{1-Lipschitz for }\rch\},\\
D(\kappa,v)&=\sup\{\Cov_{q_\kappa}(f,\langle u,v\rangle): f\ \text{1-Lipschitz for }\rho\},
\end{align*}
and $\Dch(\kappa)=\sup_{|v|=1}\Dch(\kappa,v)$, $D(\kappa)=\sup_{|v|=1}D(\kappa,v)$.
By rotation invariance about $\hat s$, both depend on $v$ only through the angle $\alpha\in[0,\pi]$ between
$v$ and $\hat s$; replacing $f$ by $-f$ shows that they bound $|\Cov|$.
\end{definition}

\begin{proposition}[Influence coefficients]\label{prop:influence}
Let $r_{\mathrm{ch}}(y,\ell)$ and $r(y,\ell)$ be Dobrushin's coefficient (5.4) for a neighbour $y$ of $\ell$, for the
metrics $\rch$ and $\rho$ respectively. Then
\[
  \beta\,\Dch(0)\;\le\;r_{\mathrm{ch}}(y,\ell)\;\le\;\beta\sup_{0\le\kappa\le6\beta}\Dch(\kappa),\qquad
  \beta\,D(0)\;\le\;r(y,\ell)\;\le\;\beta\sup_{0\le\kappa\le6\beta}D(\kappa).
\]
With $\Dch(0)=1/4$ and Theorem~\ref{thm:chordal}, for $\beta<2/9$: $r_{\mathrm{ch}}(y,\ell)=\beta/4$ and
$\sum_y r_{\mathrm{ch}}(y,\ell)=18\beta/4=9\beta/2$. With $D(0)=8/(9\pi)$ (Lemma~\ref{lem:zero}) and
Theorem~\ref{thm:onelink}, for $\beta<\pi/16$: $r(y,\ell)=8\beta/(9\pi)$ and $\sum_y r(y,\ell)=16\beta/\pi$.
\end{proposition}
\noindent Status: \st{proved} (the upper bounds use Kantorovich--Rubinstein duality \cite[Thm~10.8]{DS25},
quoted below). The geodesic upper bound is the argument of Royer \cite[(3)--(4), p.~31]{Roy86} for $\nu(h)$ on $S^n$,
here with $n=3$, $K=\beta$ and $z$ the moving staple: verbatim, ``$\mathrm R(\nu(Kz_1+\zeta),\nu(Kz_2+\zeta))
\leqslant|K|(C(n)/(n+1))^{1/2}r(z_1,z_2)$. (3)'' and, along a minimal geodesic $z(t)$,
``$\big|\frac{d}{dt}\int f(y)\nu(Kz(t)+\zeta,dy)\big|=|K\operatorname{cov}(f(y),y.z'(t))|$ (4)'' \st{read}.
The chordal upper bound is the argument of Levin, who writes the distance ``as the variation of an expectation
along an appropriate path in $\mathbf h$-space'' \cite[p.~68]{Lev80}, the straight segment
\cite[(3.35), p.~71]{Lev80} \st{read}.
Royer and Levin then bound the covariance by field-uniform constants; we instead compute $\sup_f$ of the
covariance exactly, as a function of $\kappa=|h|$.

\begin{proof}
\emph{Upper bounds.} Let $U,U'$ differ only at $y$, and let $W=AU_y^{\pm1}B$ be the unique staple
containing $y$ ($A,B$ fixed). The map $U_y\mapsto\overline{AU_y^{\pm1}B}$ is a composition of left and right
multiplications by unit quaternions and of quaternion conjugation, hence the restriction to $S^3$ of an
orthogonal map of $\R^4$; it is an isometry for both $\rch$ and $\rho$.
\emph{Chordal:} let $s(t)=s+t(s'-s)$, $t\in[0,1]$, the straight segment between the two values of $s$; then
$|s(t)|\le6$ by convexity and $|s'-s|=|\overline{W'}-\overline W|=\rch(U_y,U'_y)$. For $f$ 1-Lipschitz for $\rch$,
\[
\int f\,d(q_{s'}-q_{s})=\int_0^1\beta\,\Cov_{q_{s(t)}}\big(f,\langle u,s'-s\rangle\big)\,dt
\le\beta\,\rch(U_y,U'_y)\sup_{\kappa\le6\beta}\Dch(\kappa).
\]
\emph{Geodesic:} along a constant-speed geodesic $t\mapsto U_y(t)$, $t\in[0,1]$, from
$U_y$ to $U'_y$, the vector $s(t)$ moves on a great circle of the unit sphere translated by the
other five staples, with $|\dot s(t)|=\rho(U_y,U'_y)$, and $|s(t)|\le6$. For $f$ 1-Lipschitz for $\rho$,
\[
\int f\,d(q_{s(1)}-q_{s(0)})=\int_0^1\beta\,\Cov_{q_{s(t)}}\big(f,\langle u,\dot s(t)\rangle\big)\,dt
\le\beta\,\rho(U_y,U'_y)\sup_{\kappa\le6\beta}D(\kappa).
\]
In both cases the supremum over 1-Lipschitz $f$ of the left side is Dobrushin's Wasserstein distance (3.1)
for the corresponding metric, by the Kantorovich--Rubinstein duality as stated in \cite[Thm~10.8, p.~53]{DS25} (verbatim):
``Let $\mu$ and $\nu$ be two probability measures on a metric space $(\mathcal X,\rho)$ such that
$\int_{\mathcal X}\rho(x,x_0)\,\mu(dx)<+\infty$ and $\int_{\mathcal X}\rho(x,x_0)\,\nu(dx)<+\infty$, for some
$x_0\in\mathcal X$. Then the Wasserstein metric equals the Kantorovich-Rubinstein metric, i.e.,
$\rho_W(\mu,\nu)=\sup_{f\in\mathrm{Lip}_1(\mathcal X)}\big(\int_{\mathcal X}f(x)\,\mu(dx)-\int_{\mathcal X}f(x)\,\nu(dx)\big)$''
\st{read}, where $\rho_W$ is the infimum over couplings \cite[(2.1), p.~6]{DS25}, i.e.\ Dobrushin's (3.1).
The moment condition is trivial on the compact $S^3$. Dividing by the distance of $U_y$ and $U'_y$ and taking the
supremum gives the bounds on (5.4).

\emph{Lower bounds.} The $18$ neighbours are distinct and each staple involves three of them only, so the
six staples can be prescribed independently. Given a unit $v\in\R^4$, choose the staple containing $y$ so
that $w:=\overline{W}\perp v$, and the other five staples (unit vectors) with sum $-w$ (possible since
$|{-w}|=1\le5$); then $s=0$. Moving $U_y$ along a unit-speed geodesic moves $\overline W$ isometrically along
a geodesic of $S^3$ through $w$; choosing it with initial velocity $v$ (tangent at $w$ since $v\perp w$)
gives $\overline W(\varepsilon)=\cos\varepsilon\,w+\sin\varepsilon\,v$, so $s(\varepsilon)=\varepsilon v+O(\varepsilon^2)$,
$\rho(U_y,U_y(\varepsilon))=\varepsilon$ and $\rch(U_y,U_y(\varepsilon))=2\sin(\varepsilon/2)=\varepsilon+O(\varepsilon^3)$.
For fixed 1-Lipschitz $f$ (for either metric),
$W_1(q_0,q_{s(\varepsilon)})\ge\int f\,d(q_{s(\varepsilon)}-q_0)=\beta\varepsilon\Cov_{q_0}(f,\langle u,v\rangle)
+O(\varepsilon^2)$ (only the trivial direction of the duality is used). Dividing by the distance and
letting $\varepsilon\to0$, then optimising $f$, gives $r_{\mathrm{ch}}(y,\ell)\ge\beta\Dch(0,v)$ and
$r(y,\ell)\ge\beta D(0,v)$; both zero-field values are independent of $v$ (Lemmas~\ref{lem:refl}
and~\ref{lem:zero}).
\end{proof}

\section{The chordal one-link identity}\label{sec:chordal}

Throughout this section $\hat s=e_1$, $x_1=\langle u,e_1\rangle$, $x_2=\langle u,e_2\rangle$ with $e_2\perp e_1$,
$\bar\sigma$ is the normalised Haar measure of $S^3$, and $dq_\kappa=e^{\kappa x_1}\,d\bar\sigma/\Zb(\kappa)$ with
\[
\Zb(\kappa)=\E_{\bar\sigma}e^{\kappa x_1}=\sum_{k\ge0}\frac{\kappa^{2k}}{4^k\,k!\,(k+1)!}=\frac{2I_1(\kappa)}{\kappa}
\]
(the marginal law of $x_1$ under $\bar\sigma$ is $\frac2\pi\sqrt{1-t^2}\,dt$). We write $\E_\kappa$ for $\E_{q_\kappa}$
and $m_1=\E_\kappa x_1=\Zb'/\Zb$.

\begin{theorem}[Chordal one-link identity]\label{thm:chordal}
For all $\kappa\in[0,4/3]$,
\[
\Dch(\kappa)=\Dch(\kappa,e_2)=\E_\kappa[x_2^2]=\frac13\Big(1-\frac{\Zb''(\kappa)}{\Zb(\kappa)}\Big),
\]
\[
\frac14-\Dch(\kappa)=\frac{4\Zb''(\kappa)-\Zb(\kappa)}{12\,\Zb(\kappa)}\;\ge\;\frac{\kappa^2}{96\,\Zb(\kappa)}\;\ge\;0.0084\,\kappa^2 .
\]
In particular $\Dch(\kappa)\le1/4$, with equality only at $\kappa=0$, and
$\Dch(\kappa)=\frac14-\frac{\kappa^2}{96}+\frac{\kappa^4}{1536}-\frac{\kappa^6}{23040}+\frac{13\kappa^8}{4423680}-\cdots$.
Since $6\cdot2/9=4/3$, this covers $\kappa\le6\beta$ for all $\beta<2/9$.
\end{theorem}
\noindent Status: $\kappa=0$ and the lower bound $\Dch(\kappa)\ge\E_\kappa[x_2^2]$ \st{proved} (Lemma~\ref{lem:refl});
the upper bound \st{proved} (\S\S\ref{sec:chfield}--\ref{sec:chalpha}) up to the one-variable inequality
\eqref{eq:1d}, which is \st{interval-certified} (\S\ref{sec:chcap}; its last rational step is \st{Lean-checked});
the identity for $\frac14-\Dch$ and the bound $\kappa^2/(96\Zb)$ \st{proved}; the constant $0.0084$
\st{interval-certified}; the expansion \st{proved} (exact rational arithmetic). The inequality $6\beta<4/3$ for
$\beta<2/9$ is \st{Lean-checked} (\texttt{kappa\_max\_lt}).

\subsection{Reflection: the value at zero field and the lower bound}\label{sec:chrefl}

\begin{lemma}[Reflection]\label{lem:refl}
Let $v$ be a unit vector, $R_v x=x-2\langle x,v\rangle v$, and $q$ a probability measure on $S^3$ with
$R_{v\#}q=q$. Then $\sup\{\Cov_q(f,\langle x,v\rangle): f \text{ 1-Lipschitz for }\rch\}=\E_q\langle x,v\rangle^2$.
In particular $\Dch(0,v)=\E_{\bar\sigma}\langle x,v\rangle^2=1/4$ for every $v$, and
$\Dch(\kappa,e_2)=\E_\kappa[x_2^2]$ for every $\kappa\ge0$. \st{proved}
\end{lemma}
\begin{proof}
Put $u=\langle x,v\rangle$; $u\circ R_v=-u$, so $\E_qu=0$ and $\Cov_q(f,u)=\E_q[fu]=\frac12\E_q[(f(x)-f(R_vx))u]$.
Since $|x-R_vx|=2|u|$, this is $\le\E_q u^2$, with equality for $f=u$, which is 1-Lipschitz. For $q=\bar\sigma$,
$\E u^2=\frac14\E|x|^2=\frac14$; for $q=q_\kappa$ and $v=e_2$, $R_{e_2}$ preserves $x_1$, hence $q_\kappa$.
\end{proof}
The case $v\perp\mathbf h$ of Lemma~\ref{lem:refl} is Levin's Lemma~3.2: for $\hat{\mathbf e}\perp\mathbf h$,
``$|\mu_{\mathbf h}[(\mathbf x\cdot\hat{\mathbf e})f]|\leqslant\frac1N$ (3.12)'', proved by the same reflection
``$\leqslant\mu_{\mathbf h}(x_1^2)$ (3.14)'' \cite[p.~68]{Lev80} \st{read}. The content of
Theorem~\ref{thm:chordal} is that, for $N=4$ and $\kappa\le4/3$, no other direction does better.

\subsection{Divergence-free fields}\label{sec:chfield}

\begin{lemma}[Flux bound]\label{lem:flux}
Let $u=\langle x,v\rangle$ and $w=(u-\E_\kappa u)\,e^{\kappa x_1}/\Zb(\kappa)$ on $S^3$, so that
$\Cov_\kappa(f,u)=\E_{\bar\sigma}[fw]$. Let $V\in C^1(\bar B;\R^4)$ on the closed unit ball $\bar B\subset\R^4$ with
$\operatorname{div}V=0$ in $B$ and $V(x)\cdot x=w(x)$ on $S^3$. Then, for every $f$ 1-Lipschitz for $\rch$,
\[
\Cov_\kappa(f,u)\le\frac1{2\pi^2}\int_B|V(y)|\,dy .
\]
\st{proved}
\end{lemma}
\begin{proof}
McShane's extension $F(y)=\min_{x\in S^3}(f(x)+|y-x|)$ is 1-Lipschitz on $\R^4$ and equals $f$ on $S^3$ (here the
chordal metric is essential: it is the restriction of the Euclidean one); by Rademacher's theorem
$|\nabla F|\le1$ a.e. Since $FV$ is Lipschitz on $\bar B$ and $\operatorname{div}(FV)=\nabla F\cdot V$ a.e., the
divergence theorem and $|S^3|=2\pi^2$ give
$\E_{\bar\sigma}[fw]=\frac1{2\pi^2}\int_{S^3}F\,V\cdot x\,d\sigma=\frac1{2\pi^2}\int_B\nabla F\cdot V\,dy
\le\frac1{2\pi^2}\int_B|V|\,dy$.
\end{proof}
For $\kappa=0$ the constant field $V=v$ is admissible and gives $|B|/|S^3|=(\pi^2/2)/(2\pi^2)=1/4$: a second proof
of $\Dch(0)\le1/4$. Applied to $F(y)=\langle v,y\rangle$ (equality case), the same computation gives
$\frac1{2\pi^2}\int_B\langle v,V\rangle\,dy=\E_{\bar\sigma}[uw]=\Var_\kappa(u)$ for every admissible $V$.

\paragraph{The explicit field.} Write $y=(y_1,z)$ with $z=(0,y_2,y_3,y_4)$, $R=\sqrt{1-y_1^2}$,
$E(t)=e^{\kappa t}/\Zb(\kappa)$, and let $A$ be the bounded solution on $[-1,1]$ of
\begin{equation}\label{eq:ode}
tA(t)-\tfrac13(1-t^2)A'(t)=q(t):=(t-m_1)\,e^{\kappa t}/\Zb(\kappa),\qquad
A(t)=\frac{3}{(1-t^2)^{3/2}}\int_t^1q(s)\sqrt{1-s^2}\,ds
\end{equation}
(for $t<0$ equivalently $-\int_{-1}^t$, since $\int_{-1}^1q\sqrt{1-s^2}\,ds\propto\E_\kappa x_1-m_1=0$). For
$v=\cos\alpha\,e_1+\sin\alpha\,e_2$ put
\begin{equation}\label{eq:field}
V=\cos\alpha\,\big[A(y_1)\,e_1-\tfrac13A'(y_1)\,z\big]+\sin\alpha\,E(y_1)\,e_2 .
\end{equation}
Then $\operatorname{div}V=\cos\alpha\,(A'-A')+\sin\alpha\,\partial_{y_2}E(y_1)=0$ ($z$ has three components), and on
$S^3$, where $|z|^2=1-x_1^2$, $V\cdot x=\cos\alpha\,q(x_1)+\sin\alpha\,x_2E(x_1)=(u-\cos\alpha\,m_1)e^{\kappa x_1}/\Zb=w$,
because $\E_\kappa u=\cos\alpha\,m_1$. \st{proved}
\emph{Series and regularity.} $A=\sum_{j\ge0}\kappa^jA_j$ where $A_j$ is the unique polynomial solution of
$L[A_j]=q_j$, $L[p]=tp-\frac13(1-t^2)p'$, $q_j$ the coefficient of $\kappa^j$ in $q$: $L$ maps degree $d$ to
degree $d+1$ with leading factor $1+d/3$, the homogeneous solutions $c(1-t^2)^{-3/2}$ are unbounded, and the
solvability condition $\int q_j\sqrt{1-t^2}=0$ holds order by order; so $A_j$ has degree $j$
($A_0=1$, $A_1=\frac34t$, $A_2=\frac3{10}t^2-\frac7{40}$). $A$ is holomorphic in $\kappa$ for $|\kappa|\le2$ (where
$|\Zb|\ge2-\Zb(2)>0.4$) and bounded there, so Cauchy's estimates and Markov's inequality $\|p'\|_\infty\le
(\deg p)^2\|p\|_\infty$ make the series and its first two $t$-derivatives converge uniformly on $[-1,1]$ for
$\kappa\le4/3$; hence $V\in C^1(\bar B)$ \st{proved}. The explicit bounds are those of the certificate
(\S\ref{sec:chcap}).

\subsection{Elimination of the direction}\label{sec:chalpha}

\begin{proposition}\label{prop:alpha}
Let $m(y_1)=\min(A,E)-|A'|R/6$ and assume $m>0$ on $[-1,1]$. With
\[
J_0=\frac{2}{5\pi}\int_{-1}^1\Big(\frac{A'}{3}\Big)^2\frac{R^5}{2m}\,dy_1,\qquad
J_1'=\frac{2}{3\pi}\int_{-1}^1(E-A)^2\frac{R^3}{2m}\,dy_1,\qquad J_1=J_1'+\tfrac23J_0,
\]
one has, for all $\alpha$,
$\Dch(\kappa,\alpha)\le\cos^2\alpha\,\big(\Var_\kappa x_1+\max(J_0,J_1)\big)+\sin^2\alpha\,\E_\kappa x_2^2$.
\st{proved}
\end{proposition}
\begin{proof}
By Lemma~\ref{lem:flux} applied to \eqref{eq:field}, $\Dch(\kappa,\alpha)\le\frac1{2\pi^2}\int_B|V|$.
Pointwise, $\langle v,V\rangle=\cos^2\!\alpha\,A+\sin^2\!\alpha\,E-\cos\alpha\sin\alpha\,\frac{A'}3y_2\ge m(y_1)>0$ (as
$|\cos\alpha\sin\alpha|\le\frac12$ and $|y_2|\le R$), and $\sqrt{a^2+b^2}\le a+b^2/(2a)$ for $a>0$ gives
$|V|\le\langle v,V\rangle+|V_\perp|^2/(2m)$, where, in the orthonormal frame $(v,n,e_3,e_4)$ with
$n=-\sin\alpha\,e_1+\cos\alpha\,e_2$,
$|V_\perp|^2=\cos^2\!\alpha\,\big[\big(\sin\alpha(E-A)-\cos\alpha\frac{A'}3y_2\big)^2+\big(\frac{A'}3\big)^2(y_3^2+y_4^2)\big]$.
The integral of $\langle v,V\rangle$ is $2\pi^2\Var_\kappa(u)=2\pi^2(\cos^2\!\alpha\Var_\kappa x_1+\sin^2\!\alpha\,\E_\kappa x_2^2)$
(end of \S\ref{sec:chfield}; $\Cov_\kappa(x_1,x_2)=0$ and $\E_\kappa x_2=0$ by the symmetry $x_2\mapsto-x_2$).
On the section $\{|z|\le R\}$ of the ball, the cross term is odd in $y_2$, and $\int1=\frac43\pi R^3$,
$\int y_2^2=\frac4{15}\pi R^5$, $\int(y_3^2+y_4^2)=\frac8{15}\pi R^5$. Dividing by $2\pi^2$, with $s=\sin^2\alpha$,
the excess is at most $\cos^2\!\alpha\,\big[sJ_1'+(1-\frac s3)J_0\big]$, affine in $s\in[0,1]$, hence at most
$\cos^2\!\alpha\,\max(J_0,J_1)$.
\end{proof}

\paragraph{Reduction to one variable.} If, for a given $\kappa$, $m>0$ on $[-1,1]$ and
$\Var_\kappa x_1+\max(J_0,J_1)\le\E_\kappa x_2^2$, Proposition~\ref{prop:alpha} gives $\Dch(\kappa,\alpha)\le\E_\kappa x_2^2$
for all $\alpha$ (the convex combination \texttt{etape\_alpha}, \st{Lean-checked}), and Lemma~\ref{lem:refl} gives
equality at $\alpha=\pi/2$. Now $N:=\Zb^2(\E_\kappa x_2^2-\Var_\kappa x_1)=\frac13\Zb^2-\frac43\Zb''\Zb+\Zb'^2=\kappa^2\tilde N(\kappa)$ with
$\tilde N$ a power series, and $A'=\kappa\tilde A$, $E-A=\kappa\tilde D$ with power series $\tilde A,\tilde D$ (since
$A_0=E_0=1$). Writing $\tilde J_0,\tilde J_1$ for $J_0,J_1$ with $A',E-A$ replaced by $\tilde A,\tilde D$ (so
$J_i=\kappa^2\tilde J_i$), it suffices to prove
\begin{equation}\label{eq:1d}
\frac{\tilde N(\kappa)}{\Zb(\kappa)^2}>\max\big(\tilde J_0(\kappa),\tilde J_1(\kappa)\big)\quad\text{and}\quad m>0\ \text{on }[-1,1],
\qquad\text{for all }\kappa\in[0,4/3]
\end{equation}
(the step from \eqref{eq:1d} to the inequality before it is \texttt{etape\_kappa2}, \st{Lean-checked}). At $\kappa=0$,
exactly, $\tilde N/\Zb^2=1/48$, $\tilde A=\frac34$, $\tilde D=\frac t4$, $m=1$, and $\tilde J_0=\tilde J_1=1/256$: the
margin is $13/768$, so \eqref{eq:1d} is \emph{not} degenerate at zero field, unlike the geodesic inequality of
\S\ref{sec:onelink}. \st{proved}

\paragraph{The bound by $1/4$.} $\E_\kappa x_2^2=\frac13(1-\E_\kappa x_1^2)$ and $\E_\kappa x_1^2=\Zb''/\Zb$, so
$\frac14-\E_\kappa x_2^2=(4\Zb''-\Zb)/(12\Zb)$. The coefficient of $\kappa^{2k}$ in $4\Zb''-\Zb$ is
$\frac{3k}{k+2}\cdot\frac{1}{4^kk!(k+1)!}\ge0$ (from $z_{2k+2}/z_{2k}=\frac1{4(k+1)(k+2)}$), equal to $\frac18$ for
$k=1$. Hence $\frac14-\E_\kappa x_2^2\ge\kappa^2/(96\Zb(\kappa))\ge\kappa^2/(96\Zb(4/3))$, and
$\Zb(4/3)\in[1.2393065416564178344\pm3\cdot10^{-20}]$ gives $1/(96\Zb(4/3))\ge0.0084052$ \st{interval-certified}.
The expansion of Theorem~\ref{thm:chordal} follows by exact series division. \st{proved}

\subsection{Computer-assisted part}\label{sec:chcap}

\paragraph{Interval evaluation.} Script \path{certificate_chordal/cert_ch.py}, Arb ball arithmetic
(python-flint 0.9.0), working precision $128$ bits and series length cap $200$, both mandatory arguments, set
explicitly and read back. All series in $\kappa$ ($\Zb,\Zb',\Zb'',1/\Zb$, $m_1$, $m_1/\Zb$, $\Zb^2$, $\Zb''\Zb$, $\Zb'^2$,
$N$) are computed to length $L=D+5=85$ ($D=80$); guard~1 refuses $\mathrm{cap}<L$, guard~2 asserts
\texttt{.prec}${}=L$ on each of these ten series and products (trailing zeros are padded explicitly, since
\texttt{coeffs()} drops them). The $A_j$ are obtained by exact downward recursion; the solvability residual
contains $0$ for all $j\le80$. \emph{Tails:} all orders $>80$ are bounded, not neglected: Cauchy's estimates on
$|\kappa|\le2$ ($|\Zb|\ge0.409363$, $|m_1|\le1.6830$, $|E|\le18.050$, $|A|\le228.21$, using
$\int_t^1\sqrt{1-s^2}\,ds\le\frac\pi2(1-t^2)^{3/2}$) and on $|\kappa|\le4$ for $N$, and Markov's inequality for
$A_j'$. At $\kappa=4/3$ the tails are $1.9\cdot10^{-8}$ ($\tilde A$), $3.0\cdot10^{-12}$ ($\tilde D$) and
$7.3\cdot10^{-39}$ ($N$).
\emph{Boxes:} $[0,4/3]$ is cut into $256$ boxes $[i/192,(i+1)/192]$, $\kappa$ being a ball in each; the series are
folded into polynomials in $t=y_1$, and $[-1,1]$ is cut into $1024$ dyadic sub-intervals, on which the polynomials
are evaluated in centred form $p(t_c)+p'(T)(T-t_c)$. On each sub-interval $m>0$ is \emph{asserted}, and $\tilde J_0$,
$\tilde J_1$ are bounded above by upper Riemann sums (with exact upper bounds for $R^3,R^5$). $\tilde N$ is bounded
below on the box (minus its tail), and $\Zb$ above by its value at the right end (nonnegative coefficients).

\paragraph{Result.} All $256$ boxes pass; the worst margin is $1.142\cdot10^{-2}$, on the last box
$[85/64,4/3]$ (lower bound $1.6883\cdot10^{-2}$ for $\tilde N/\Zb^2$ against $\tilde J_1\le5.463\cdot10^{-3}$,
$\tilde J_0\le3.975\cdot10^{-3}$, $m\ge0.192$); at $\kappa\in[0,1/192]$ the margin is $1.690\cdot10^{-2}$.
The right-hand side can be multiplied by about $3$ before the inequality fails. \st{interval-certified}

\paragraph{Gates.} All logs return code 0 unless a refusal is expected (logs in the archive).
\begin{center}\small
\begin{tabular}{p{0.30\linewidth}p{0.62\linewidth}}
\toprule
gate & outcome\\
\midrule
cap $\in\{85,100,200,300\}$ & identical results in all $256$ boxes; cap $10$ and $84={}L-1$ refused by guard~1, and by guard~2 when guard~1 is disabled\\
order $D\in\{40,60,80,100\}$ & the bound moves ($\tilde J_0$ of the last box $4.8604\cdot10^{-3}$, $3.97576\cdot10^{-3}$, $3.975228508\cdot10^{-3}$, $3.975228233\cdot10^{-3}$) and stabilises; all four runs pass\\
sub-intervals $\in\{256,\dots,2048\}$ & upper Riemann sums strictly decreasing, step halving ($O(1/N_T)$)\\
precision $128$ vs $192$ & $\max|\Delta\tilde J_1|=6.7\cdot10^{-14}$; all boxes pass\\
boxes $\in\{64,128,512\}$ & all pass; worst margin $1.106$, $1.130$, $1.148\cdot10^{-2}$\\
exact gate (fractions, no Arb) & solvability residuals exactly $0$ for $j\le80$; $\tilde N_0=1/48$, odd orders $0$; all $6891$ exact coefficients lie in the Arb balls\\
numerical gate (\st{numerical}) & series $A$ vs direct quadrature $1.3\cdot10^{-15}$; $|\operatorname{div}V|\le4\cdot10^{-11}$; $|V\cdot x-w|\le7\cdot10^{-16}$; full flux bound $\frac1{2\pi^2}\int_B|V|$ by 3D quadrature within $1.7\cdot10^{-4}$ of the linear programs of \path{calc/lp2d_m14.log} (e.g.\ $0.20892$, $0.22140$, $0.23332$ vs $0.20887$, $0.22145$, $0.23324$ at $\kappa=4/3$, $\alpha=0,\pi/4,\pi/2$)\\
\bottomrule
\end{tabular}
\end{center}

\paragraph{Independent check.} A second implementation (\path{independent_check_chordal/}), written without
reading or importing the certificate code, recomputes everything with different tools: exact rational
coefficients up to order $100$ (own downward solver; solvability residuals exactly $0$ and $L[A_j]=q_j$ checked
exactly), \texttt{mpmath.iv} interval arithmetic instead of Arb, elementary tail bounds instead of complex Cauchy
estimates ($|[\kappa^n]1/\Zb|\le2^{-n}$ by induction, and $(1-t^2)^{-3/2}\int_t^1\sqrt{1-s^2}\,ds\le1$ on $[0,1)$ instead of
$\pi/2$), a uniform partition in $\theta$ ($t=\cos\theta$) with exact weights $\int\sin^4$, $\int\sin^6$, and naive
Horner evaluation. On its own partition ($200$ boxes of width $1/150$) it certifies \eqref{eq:1d} on all of
$[0,4/3]$ (worst margin $1.1395\cdot10^{-2}$, $m\ge0.188$); it also recertifies the $256$ boxes of the Lean data
file and checks, two-sidedly, that none of the exported rationals contradicts its own enclosures. Its negative
gates fail as they should (precision sabotaged to $53$ bits: refused; series shorter than $D+2$: refused; left
side of \eqref{eq:1d} scaled by $0.30$: fails on the last box, while $0.35$ passes).
Exact checks: $\E_{\bar\sigma}u^2=1/4$ and the cost $1/4$ of the reflection coupling (SymPy), the coefficient
$3k/(k+2)$ symbolically for all $k$, $c_2=1/96$; a Monte-Carlo evaluation of the full flux bound at $\kappa=4/3$
($0.20884$, $0.22131$, $0.23322$) agrees with the linear programs \st{numerical}. \st{interval-certified}

\paragraph{Lean.} Commit \path{1b0f039c53c1e2a11c496cc5a1cb7774ab5a010e} of the author's Lean library (parent
\texttt{7a211ed}, which holds the geodesic modules of \S\ref{sec:cap}; three new files, no existing file changed),
module \path{KevYM/Certif/DobrushinCordale/}. The three modules are reproduced byte for byte in the standalone
project \path{lean/} of the accompanying archive. Lean toolchain \texttt{v4.34.0}. Main statements (verbatim):
{\scriptsize
\begin{verbatim}
theorem sound (L : List Boite) (a c : ℚ) (hch : chaine L a c = true) (hok : L.all Boite.ok = true)
    (Nt J0 J1 : ℝ → ℝ)
    (hN : ∀ b ∈ L, ∀ κ : ℝ, (b.k0 : ℝ) ≤ κ → κ ≤ (b.k1 : ℝ) → (b.lb : ℝ) ≤ Nt κ)
    (hJ0 : ∀ b ∈ L, ∀ κ : ℝ, (b.k0 : ℝ) ≤ κ → κ ≤ (b.k1 : ℝ) → J0 κ ≤ (b.h0 : ℝ))
    (hJ1 : ∀ b ∈ L, ∀ κ : ℝ, (b.k0 : ℝ) ≤ κ → κ ≤ (b.k1 : ℝ) → J1 κ ≤ (b.h1 : ℝ))
    (κ : ℝ) (ha : (a : ℝ) ≤ κ) (hc : κ ≤ (c : ℝ)) : max (J0 κ) (J1 κ) < Nt κ
theorem etape_kappa2 (k Nt Jt V1 E2 : ℝ) (hE : E2 - V1 = k ^ 2 * Nt) (h : Jt < Nt) :
    V1 + k ^ 2 * Jt ≤ E2
theorem etape_alpha (V1 E2 J c s : ℝ) (hcs : c ^ 2 + s ^ 2 = 1) (h : V1 + J ≤ E2) :
    c ^ 2 * (V1 + J) + s ^ 2 * E2 ≤ E2
theorem somme_dobrushin_lt_one (β : ℝ) (hβ : β < 2 / 9) : 18 * β * (1 / 4) < 1
theorem kappa_max_lt (β : ℝ) (hβ : β < 2 / 9) : 6 * β < 4 / 3
theorem pi_div_16_lt : Real.pi / 16 < 2 / 9
theorem certificat_cordal (Nt J0 J1 : ℝ → ℝ)
    (hN : ∀ b ∈ boites, ∀ κ : ℝ, (b.k0 : ℝ) ≤ κ → κ ≤ (b.k1 : ℝ) → (b.lb : ℝ) ≤ Nt κ)
    (hJ0 : ∀ b ∈ boites, ∀ κ : ℝ, (b.k0 : ℝ) ≤ κ → κ ≤ (b.k1 : ℝ) → J0 κ ≤ (b.h0 : ℝ))
    (hJ1 : ∀ b ∈ boites, ∀ κ : ℝ, (b.k0 : ℝ) ≤ κ → κ ≤ (b.k1 : ℝ) → J1 κ ≤ (b.h1 : ℝ))
    (κ : ℝ) (h0 : 0 ≤ κ) (hK : κ ≤ 4 / 3) : max (J0 κ) (J1 κ) < Nt κ
\end{verbatim}}
Here \texttt{Boite.ok b} is \texttt{h0 < lb \&\& h1 < lb \&\& k0 ≤ k1}, \texttt{chaine} checks that the boxes are
chained from $a$ to $c$, and \texttt{boites} is the list of the $256$ exported boxes $\langle k_0,k_1,lb,h_0,h_1\rangle$
(Arb bounds rounded outward to the grid $2^{-40}$, re-checked exactly before export); \texttt{boites\_ok} and
\texttt{boites\_chaine} are proved by \texttt{decide +kernel}. The file also proves $18\cdot(2/9)\cdot\frac14=1$, the
converse for $\beta\ge2/9$, and $6\cdot(2/9)=4/3$. Build clean (0 warnings); all $16$ theorems of the module depend
only on \texttt{propext}, \texttt{Classical.choice}, \texttt{Quot.sound}; the rebuild from \texttt{git archive 1b0f039}
succeeds, and the standalone project \path{lean/} of the archive builds from an empty \texttt{.lake} with the same
axiom report (\path{lean/logs/build_from_zero.log}). \st{Lean-checked}

\paragraph{What Lean covers, exactly.} \texttt{certificat\_cordal} is a conditional implication: \emph{if} real
functions \texttt{Nt}, \texttt{J0}, \texttt{J1} (standing for $\tilde N/\Zb^2$, $\tilde J_0$, $\tilde J_1$) are bounded box by
box by the exported rationals, \emph{then} $\max(\tilde J_0,\tilde J_1)<\tilde N/\Zb^2$ on $[0,4/3]$. Lean also checks
the two algebraic steps \texttt{etape\_alpha}, \texttt{etape\_kappa2} and the threshold arithmetic. Not in Lean:
the flux bound, the explicit field and Proposition~\ref{prop:alpha} (proved on paper, \S\S\ref{sec:chfield}--\ref{sec:chalpha}),
the bounds of each box and the positivity of $m$ (\st{interval-certified}), and the constant $0.0084$.

\paragraph{Artefacts (SHA-256).} Paths are relative to the accompanying archive.\par\noindent
{\small
\begin{tabularx}{\linewidth}{@{}>{\raggedright\arraybackslash}p{0.40\linewidth}>{\raggedright\arraybackslash}X@{}}
\toprule
file & SHA-256\\
\midrule
\path{certificate_chordal/cert_ch.py} & \path{d83aa23b2fd20707b08076c87053aa8c9bc9e277ca15d78be95e5afa7dab14c0}\\
\path{certificate_chordal/export_lean.py} & \path{f9f59cbeb0a0ac0c609ce7471b97ea3c69ea7f7330ac4c5975c152416b374ac4}\\
\path{certificate_chordal/cert_D80_NT1024_NK256_cap200_p128.json} & \path{aa5da0f05c1ee230782c30bd451981df2100ca75a9fc46ca73a6b7e54c765bb5}\\
\path{lean/KevYM/Certif/DobrushinCordale/Donnees.lean} (${}={}$\texttt{1b0f039}) & \path{72a0901472475b3a6d98033692b899e32c5f9b2a1daab7b9ed5c22416dc462fe}\\
\bottomrule
\end{tabularx}}

\section{The geodesic one-link bound}\label{sec:onelink}

This section is the main result of v1.2 of this note, unchanged in substance. It gives the exact
field-uniform coefficient for the geodesic metric, hence Theorem~\ref{thm:geo}.

\begin{theorem}[One-link bound, geodesic metric]\label{thm:onelink}
For all $\kappa\in[0,59/50]$ and all unit $v$,
\[
D(\kappa,v)\;\le\;\frac{8}{9\pi}-0.0085\,\kappa^2 ,
\]
with equality at $\kappa=0$. Since $6\cdot\pi/16=3\pi/8<59/50$, this covers $\kappa\le6\beta$ for all
$\beta<\pi/16$.
\end{theorem}
\noindent Status: $\kappa=0$ \st{proved}; $\kappa>0$ \st{interval-certified}. Precisely: Lean proves only
that the monotone bound \texttt{bP} is $<0$ in each of the $160$ boxes, i.e.\ $\Phi<0$ on $[0,K^2]$
\emph{if} the exported rationals bound the coefficients (\st{Lean-checked}, \S\ref{sec:cap}); the
constant $0.0085$ and the covering of $\alpha\in[0,\pi]$ by the boxes (box $\mapsto$ $\alpha$-interval) are
\st{interval-certified}, not in Lean. The inequality $3\pi/8<59/50$ is \st{Lean-checked}
(\texttt{kappa\_max\_lt}).

\begin{lemma}[$\kappa=0$]\label{lem:zero}
$D(0,v)=8/(9\pi)$ for every unit $v$. \st{proved}
\end{lemma}
\begin{proof}
Let $\gamma\in[0,\pi]$ be the angle between $u$ and $v$; under Haar measure its law is
$\frac2\pi\sin^2\gamma\,d\gamma$. \emph{Upper bound:} $\cos\gamma$ is a first eigenfunction of
$-\Delta_{S^3}$ (eigenvalue $3$), so $\psi=\cos\gamma/3$ solves $-\Delta\psi=\langle u,v\rangle$ and, for
1-Lipschitz $f$, $\Cov(f,\langle u,v\rangle)=\E[\nabla f\cdot\nabla\psi]\le\E|\nabla\psi|
=\frac13\cdot\frac2\pi\int_0^\pi\sin^3\gamma\,d\gamma=\frac{8}{9\pi}$.
\emph{Lower bound:} $f=\pi/2-\gamma$ is 1-Lipschitz and (integration by parts)
\[
\Cov(f,\cos\gamma)=-\E[\gamma\cos\gamma]=-\frac2\pi\int_0^\pi\gamma\cos\gamma\sin^2\gamma\,d\gamma
=-\frac2\pi\Big(-\frac49\Big)=\frac{8}{9\pi}.
\]
\end{proof}

\subsection{Reduction (proved, not formalised)}\label{sec:reduction}

\paragraph{Step 1: transport.} Let $L=\Delta_{S^3}+\kappa\nabla x\cdot\nabla$ with $x=\langle u,\hat s\rangle$;
$L$ is symmetric in $L^2(q_\kappa)$, with Dirichlet form $\E_{q_\kappa}[\nabla\varphi\cdot\nabla\psi]$.
\emph{Weak solutions (Lax--Milgram).} Let $H$ be the space of $\varphi\in H^1(S^3)$ with $\E_{q_\kappa}\varphi=0$,
with the inner product $\E_{q_\kappa}[\nabla\varphi\cdot\nabla\psi]$; it is a Hilbert space by the Poincar\'e
inequality of Step~3 ($\lambda>0$). For $g\in L^2(q_\kappa)$ with $\E_{q_\kappa}g=0$, the functional
$\varphi\mapsto\E_{q_\kappa}[\varphi g]$ is bounded on $H$ by $\|g\|_2\lambda^{-1/2}\|\nabla\varphi\|_2$, so by
Riesz there is a unique $\psi\in H$ with $\E[\nabla\varphi\cdot\nabla\psi]=\E[\varphi g]$ for all
$\varphi\in H^1$ (constants integrate against $g$ to zero); we write $-L\psi=g$ in this weak sense, which is
all that is used below. With $g=\langle u,v\rangle-\E_{q_\kappa}\langle u,v\rangle$ and $f$ 1-Lipschitz for
$\rho$ (so $f\in H^1$ with $|\nabla f|\le1$ a.e., by Rademacher's theorem),
$\Cov_{q_\kappa}(f,\langle u,v\rangle)=\E[\nabla f\cdot\nabla\psi_v]\le\E|\nabla\psi_v|$, so
$D(\kappa,v)\le\E_{q_\kappa}|\nabla\psi_v|$. \st{proved}

\paragraph{Step 2: reduced ODEs and exact series.} Write $v=\cos\alpha\,\hat s+\sin\alpha\,e$ with
$e\perp\hat s$, $y=\langle u,e\rangle$, and seek $\psi_v=\cos\alpha\,g(x)+\sin\alpha\,y\,h(x)$, $p=g'$.
Using $\Delta_{S^3}g(x)=(1-x^2)g''-3xg'$, $\Delta_{S^3}(yh)=y[(1-x^2)h''-5xh'-3h]$,
$|\nabla x|^2=1-x^2$ and $\nabla x\cdot\nabla y=-xy$, the equation reduces to
\[
(1-x^2)(p'+\kappa p)-3xp=-(x-m),\qquad (1-x^2)(h''+\kappa h')-5xh'-(3+\kappa x)h=-1,
\]
with $m=\E_{q_\kappa}x$. Expanding $p=\sum_n\kappa^np_n$, $h=\sum_n\kappa^nh_n$, $m=\sum_n\kappa^nm_n$ gives
polynomials $p_n,h_n$ with rational coefficients, computed exactly; $p_n,h_n$ have parity $(-1)^n$ and
$m_n=0$ for $n$ even (asserted in exact arithmetic by the certificate script).
For the truncation $\psi_N$ at order $N$, the residual is
$L\psi_N+(\langle u,v\rangle-\cos\alpha\,m)=\cos\alpha\,(m_N-m)+R'$ with
$R'=\kappa^{N+1}\big[\cos\alpha(1-x^2)p_N+\sin\alpha\,y\big((1-x^2)h_N'-xh_N\big)\big]$
(identity also checked symbolically in $\R^4$ at rational points of $S^3$).
Since the full residual $R$ has $q_\kappa$-mean zero, $|\cos\alpha(m_N-m)|\le\|R'\|_\infty$, hence
$\|R\|_\infty\le2\kappa^{N+1}\rho_N$, where $\rho_N$ is the sum of $\ell^1$ norms of the coefficients of the
two bracketed polynomials ($\rho_{16}\approx8.1\cdot10^{-10}$). \st{proved}

\paragraph{Step 3: correction via Poincar\'e.} Let $-L\varphi=R$; then $\psi_v=\psi_N+\varphi$ exactly. The
density of $q_\kappa$ w.r.t.\ Haar measure has log-oscillation at most $2\kappa$ and the spectral gap of
$-\Delta_{S^3}$ is $3$, so by the Holley--Stroock perturbation argument \cite{HS87} the Poincar\'e constant
of $q_\kappa$ satisfies $\lambda\ge3e^{-2\kappa}$. (Proof: write $dq_\kappa=h\,d\sigma$ with $\sigma$ the
Haar measure, $h=e^{\kappa x}/\tilde Z$, so $e^{-\kappa}/\tilde Z\le h\le e^{\kappa}/\tilde Z$; then
$\operatorname{Var}_{q}f\le\E_q(f-\E_\sigma f)^2\le\frac{e^{\kappa}}{\tilde Z}\operatorname{Var}_\sigma f
\le\frac{e^{\kappa}}{3\tilde Z}\E_\sigma|\nabla f|^2\le\frac{e^{2\kappa}}{3}\E_q|\nabla f|^2$.) Then
$\E|\nabla\varphi|^2=\E[\varphi R]\le\|R\|_2(\E|\nabla\varphi|^2/\lambda)^{1/2}$, whence
\[
\E|\nabla\varphi|\le\|R\|_2/\sqrt\lambda\le\varepsilon_N(\kappa):=2\kappa^{N+1}\rho_N e^{\kappa}/\sqrt3 .
\]
So $D(\kappa,v)\le\E|\nabla\psi_N|+\varepsilon_N(\kappa)$. \st{proved}

\paragraph{Step 4: AM--GM.} With $\varrho=\sin\gamma/3=|\nabla\psi_0|$ ($\cos\gamma=\langle u,v\rangle$),
pointwise $|X|\le(|X|^2/\varrho+\varrho)/2$, hence
$\E|\nabla\psi_N|\le U_N:=\E\big[(|\nabla\psi_N|^2/\varrho+\varrho)/2\big]$, with equality at $\kappa=0$. \st{proved}

\paragraph{Step 5: polynomial coordinates.} The image of the uniform measure on $S^3$ under
$u\mapsto(\langle u,\hat s\rangle,\langle u,e\rangle)$ is uniform on the unit disc, and $|\nabla\psi_N|^2$ is a
polynomial in $(x,y,\kappa)$. In the rotated frame $s=\langle u,v\rangle$, $t=\sqrt{1-s^2}\,\tau$, the singular
weight $1/\varrho$ cancels against the Jacobian and
\[
U_N(\kappa,\alpha)\,Z(\kappa)=\iint_{[-1,1]^2}\Big[\tfrac32|\nabla\psi_N|^2+\tfrac{1-s^2}{6}\Big]e^{\kappa x}\,ds\,d\tau,
\qquad Z(\kappa)=\pi\sum_{k\ge0}\frac{(\kappa^2/4)^k}{k!\,(k+1)!}=\pi\Zb(\kappa).
\]
At $\kappa=0$ the integral is $8/9$ and $Z(0)=\pi$, recovering $8/(9\pi)$. The coefficient of $\kappa^j$ is the
integral of a polynomial of degree $\le j+2$ in $(s,\tau)$: the terms odd in $\tau$ integrate to zero
(also exactly on symmetric Gauss nodes), and in the remaining terms $\sqrt{1-s^2}$ occurs only squared.
\st{proved}

\paragraph{Step 6: exact factorisation of $\kappa^2$.} Put
$\Delta(\kappa,\alpha):=\big(U_N+\varepsilon_N-\tfrac{8}{9\pi}\big)Z=\sum_jE_j(\alpha)\kappa^j+(\text{remainders})$.
$E_0=0$ exactly ($I_0=8/9$), and $E_j=0$ for $j$ odd by the parity $u\mapsto-u$, i.e.\ $(s,\tau)\mapsto(-s,-\tau)$.
Hence $\Delta=\kappa^2\Phi(w)$ with $w=\kappa^2$ and
$\Phi(w)=\sum_{k\ge1}E_{2k}w^{k-1}+(\text{Taylor remainder}+\varepsilon_NZ)/\kappa^2$; the last two terms
are increasing in $\kappa$ and are bounded at $\kappa=K=59/50$. This is what makes the comparison at
$\kappa\to0$ non-degenerate: no two touching balls are ever compared. \st{proved}

Since $Z$ is increasing, $\Phi\le-c_\ast$ on $[0,K^2]$ implies
$D(\kappa,v)\le\frac{8}{9\pi}-\frac{c_\ast}{Z(K)}\kappa^2$.

\begin{remark}[second order, exact]
An independent exact computation (SymPy, harmonic components in $\R^4$; code not included in the archive)
gives $E_2(\alpha)=-(1309\cos^2\alpha+779)/21600$, maximal at $\alpha=\pi/2$. Locally as $\kappa\to0$ the
method cannot give a coefficient better than $779/(21600\pi)\approx0.0115$ in front of $\kappa^2$; along the
route actually used ($c_\ast/Z(K)$, uniform on $[0,K]$), since $c_\ast\le-\sup_\alpha\Phi(0)=779/21600$, the
ceiling is $779/(21600\,Z(K))\approx0.0097$. The certified $0.0085$ is below both.
(Status: exact computation; the same $E_2$ and $E_4$ are reproduced in exact rationals, by a different method,
by the independent check \path{independent_check_geodesic/}.)
\end{remark}

\subsection{Computer-assisted part}\label{sec:cap}

\paragraph{Interval evaluation.} Parameters: series order $N=16$; Gauss--Legendre with $N_G=22$ nodes per
direction, exact for polynomial degree $\le2N_G-1=43$, hence for all $j\le D=41$; $\alpha\in[0,\pi]$ cut into
$160$ boxes of width $\pi/160$, $(\cos\alpha,\sin\alpha)$ enclosed by balls; Arb ball arithmetic
(python-flint 0.9.0) with working precision $128$ bits and series length cap $200$, both set explicitly and
read back. Nodes and weights are rigorous balls (\texttt{arb.legendre\_p\_root}). Orders $\kappa^0,\dots,\kappa^{41}$
are computed; the orders $\ge42$ of $e^{\kappa x}$ are covered by the Lagrange bound
$|\text{rem}|/\kappa^2\le4\kappa^{D-1}e^{\kappa}\sum_i\sup|A_i|/(D+1-i)!\le1.53\cdot10^{-21}$ at $\kappa=K$.
On each box, $\Phi(w)\le\sum_k\max(c_k\cdot0^k,c_kK^{2k})$ with $c_k$ upper bounds of $E_{2k+2}$ (the constant
term absorbing the remainder bounds).

\paragraph{Result.} All $160$ boxes pass in one piece in $w$; the worst bound is
$\Phi\le-0.03175361$ (box 75, $\alpha$ near $\pi/2$); $Z(59/50)\in[3.721046472\pm4.7\cdot10^{-10}]$, and
$0.03175361/Z(59/50)\in[0.0085335161\pm5\cdot10^{-12}]>0.0085$. A second run with $N=20$, $N_G=26$,
$240$ boxes, precision $192$ passes with worst bound $-0.03333053$ (stability check, same code).
\st{interval-certified}

\paragraph{Erratum on an earlier version (important).} A first version of this certificate (Lean commit
\texttt{1f3ace0}) was \emph{invalid}: python-flint's default \texttt{ctx.cap}${}=10$ silently truncated every
product of power series at $\kappa^9$, so the orders $\kappa^{10},\dots,\kappa^{41}$ were neither computed nor
bounded, and in each of the $160$ boxes at least one exported bound (orders $E_{12}$--$E_{40}$) was below
the true value. Lean accepted it, since Lean only sees the rational inequalities. The error was found by
an independent adversarial re-verification, and its numerical effect is small ($\le4.7\cdot10^{-6}$ against a margin
$0.0317$), but the artefact had to be regenerated. The repaired script makes cap and precision mandatory
arguments, refuses $\mathrm{cap}<L=42$ (guard~1) and asserts that every series and product has length
exactly $L$ (guard~2, $851\,840$ assertions per run; the assertion is on \texttt{.prec} because
\texttt{coeffs()} drops trailing zeros). The chordal certificate of \S\ref{sec:chcap} was written with the
same guards from the start.

\paragraph{Independent checks (gates).} All logs return code 0 unless a refusal is expected.
\begin{center}\small
\begin{tabular}{p{0.28\linewidth}p{0.64\linewidth}}
\toprule
gate & outcome\\
\midrule
G1: cap $\in\{42,50,100,200\}$ & identical results in all $160$ boxes; cap $=10$ refused by guard~1; cap $=10$ and $41$ with guard~1 disabled refused by guard~2\\
G2: explicit Cauchy convolution (no \texttt{arb\_series}) & all $42$ coefficients overlap; $\max|\Delta\mathrm{mid}|=1.5\cdot10^{-39}$\\
G3: $N\in\{16,20,24\}$ vs $32$ & $|U_N-U_{32}|$ nonzero, decreasing ($1.34\cdot10^{-10}$, $1.20\cdot10^{-12}$, $1.08\cdot10^{-14}$ at $\alpha=0$, $\kappa=1.18$), below $\varepsilon_N$\\
G4: $N_G\in\{21,22,26\}$ & $22$ vs $26$ overlap for $j\le41$; with $21$ nodes disagreement starts exactly at $j=40$, as predicted by Gauss exactness\\
G5: Arb vs float polar grid & gap $\le1.0\cdot10^{-8}$, decreasing as $O(h^2)$ under grid refinement\\
independent recomputation (Galerkin; code not included in the archive) & Galerkin solver at finite $\kappa$ and exact $E_2,E_4$: consistent with every box bound; recomputed data identical byte for byte\\
independent check (exact rationals; \path{independent_check_geodesic/}) & $E_j(\alpha)\in\mathbb Q[\cos\alpha]$ computed exactly by a different method (no quadrature, no \texttt{arb\_series}); all $3200$ Lean coefficients of \texttt{7a211ed} ($160$ boxes $\times$ $E_2,\dots,E_{40}$) are true upper bounds, via rigorous box majorants with the invariant ``majorant $\ge$ exact sample'' checked everywhere; the same tool refutes the old commit \texttt{1f3ace0} ($1634$ false bounds, $160/160$ boxes); worst \texttt{bP} $=-0.03175360$ in rationals; $0.0085335\ge0.0085$ checked exactly against a rigorous upper bound of $Z(K)$; $U_{16}$ vs.\ an independent Galerkin solver at $700$ points $(\kappa,\alpha)$: $\max|\Delta|=1.34\cdot10^{-10}<\varepsilon_{16}$\\
\bottomrule
\end{tabular}
\end{center}
Floating-point gates are \st{numerical}; they test the modelling, they are not part of the proof.

\paragraph{Lean.} Commit \path{7a211ed0426e95b36a0ff35f5e896f8594b34fd3} of the author's Lean library (parent
\texttt{1f3ace0}; only \path{KevYM/Certif/DobrushinLien/Donnees.lean} changed). The three modules are reproduced byte
for byte in the standalone Lean project \path{lean/} of the accompanying archive (they are unchanged in
\texttt{1b0f039}). Lean toolchain \texttt{v4.34.0}. Main statements (verbatim):
{\footnotesize
\begin{verbatim}
theorem sound (cs : List ℚ) (a b : ℚ) (ha : 0 ≤ a) (hneg : bP cs 0 a b < 0)
    (es : List ℝ) (hes : List.Forall₂ (fun (e : ℝ) (c : ℚ) => e ≤ (c : ℝ)) es cs)
    (w : ℝ) (haw : (a : ℝ) ≤ w) (hwb : w ≤ b) : pR es 0 w < 0
theorem certs_ok : certs.all (fun cs => decide (bP cs 0 0 KW < 0)) = true
theorem certificat_boites (cs : List ℚ) (hcs : cs ∈ certs) (es : List ℝ)
    (hes : List.Forall₂ (fun (e : ℝ) (c : ℚ) => e ≤ (c : ℝ)) es cs)
    (w : ℝ) (hw0 : 0 ≤ w) (hw : w ≤ (KW : ℝ)) : pR es 0 w < 0
theorem somme_dobrushin_lt_one (β : ℝ) (hβ : β < Real.pi / 16) :
    18 * β * (8 / (9 * Real.pi)) < 1
theorem kappa_max_lt : 6 * (Real.pi / 16) < (59 : ℝ) / 50
\end{verbatim}}
Here \texttt{bP} is the monotone upper bound $\sum_k\max(c_ka^k,c_kb^k)$, \texttt{pR es 0 w}${}=\sum_je_jw^j$
(\texttt{pR\_eq\_sum}), \texttt{KW}${}=3481/2500=(59/50)^2$, and \texttt{certs} is the list of $160$ exported
rational coefficient vectors (Arb upper bounds rounded up to the grid $2^{-40}$, re-checked exactly
before export). \texttt{certs\_ok} is proved by \texttt{decide +kernel}. The file also proves
$18\cdot(\pi/16)\cdot 8/(9\pi)=1$ and the converse for $\beta\ge\pi/16$.
Build clean (0 warnings); all $13$ theorems depend only on \texttt{propext}, \texttt{Classical.choice},
\texttt{Quot.sound}; the rebuild from \texttt{git archive 7a211ed} succeeds, and so does the from-scratch build of
the standalone project (\path{lean/logs/build_from_zero.log}). \st{Lean-checked}

\paragraph{What Lean covers, exactly.} Lean proves \texttt{bP cs 0 0 KW < 0} (strict) for each of the
$160$ exported vectors, and the implication of \texttt{certificat\_boites}: \emph{if} the true
coefficients lie below the exported rationals, \emph{then} $\Phi<0$ on $[0,K^2]$; plus the threshold
arithmetic. Not in Lean: the constant $0.0085$ (i.e.\ $-\max\texttt{bP}/Z(K)\ge0.0085$, checked exactly in
rationals by the independent check against a rigorous upper bound of $Z(K)$) and the covering of $[0,\pi]$ by the boxes
(the association box $\mapsto\alpha$-interval is a comment in the Lean file); both are
\st{interval-certified}.

\paragraph{Trust base.} The implication ``integrals $\Rightarrow$ exported bounds'' rests on Arb and on the
reduction of \S\ref{sec:reduction} (proved on paper; the classical inputs are Kantorovich--Rubinstein
duality, quoted from \cite[Thm~10.8]{DS25}, the Holley--Stroock perturbation argument, written out in
Step~3, and the weak solvability of $-L\varphi=R$, written out in Step~1), and is \emph{not} formalised. It is
cross-checked in exact rationals by the independent check (table above), which shares with the certificate only the
modelling (reduced ODEs, measure $ds\,d\tau$).

\paragraph{Artefacts (SHA-256).} Paths are relative to the accompanying archive.\par\noindent
{\small
\begin{tabularx}{\linewidth}{@{}>{\raggedright\arraybackslash}p{0.40\linewidth}>{\raggedright\arraybackslash}X@{}}
\toprule
file & SHA-256\\
\midrule
\path{certificate_geodesic/cert_lien.py} (v2) & \path{8b76b11125693b8f71b44898e959635066b3804560c2ee314092667a562ea096}\\
\path{certificate_geodesic/export_lean.py} & \path{1885f8bc2e8444bd616efa41c5c18bbfaa3fca78e6bc3e06ae197014160c1bcb}\\
\path{certificate_geodesic/cert_N16_G22_B160_cap200_p128.json} & \path{91700df0e8031e1f6ed4c3d6d9e88c42ddf36f0abeacd58c9f328738b0d9c43e}\\
\path{lean/KevYM/Certif/DobrushinLien/Donnees.lean} (${}={}$\texttt{7a211ed}) & \path{9d6c06283848a9ef38baecb222e195e9aa528009412e7c8bded6f2555321c37f}\\
\bottomrule
\end{tabularx}}

\section{Proofs of Theorems~\ref{thm:main} and~\ref{thm:geo}}
\emph{Theorem~\ref{thm:main}.} Let $\beta<2/9$, and use the chordal metric. By Theorem~\ref{thm:chordal},
$\sup_{\kappa\le6\beta}\Dch(\kappa)=\Dch(0)=1/4$ (as $6\beta<4/3$, \texttt{kappa\_max\_lt}); by
Proposition~\ref{prop:influence} each of the $18$ neighbours has influence $\beta/4$ and all other links have
influence $0$; so (5.2) holds with $d=9\beta/2<1$ (\texttt{somme\_dobrushin\_lt\_one}). (5.1) holds for every
field. By \cite[Thm~4]{Dob70} there is at most one field with the Wilson single-link conditionals, hence at most
one Gibbs measure (item~2 of \S\ref{sec:dobrushin}); existence is Proposition~\ref{prop:exist}. \qed

\emph{Theorem~\ref{thm:geo}.} Let $\beta<\pi/16$, and use the geodesic metric. By Theorem~\ref{thm:onelink},
$\sup_{\kappa\le6\beta}D(\kappa)=D(0)=8/(9\pi)$, and Proposition~\ref{prop:influence} gives the sum $16\beta/\pi<1$
(\texttt{somme\_dobrushin\_lt\_one} of the geodesic module). For $\beta\ge\pi/16$ the lower bound of
Proposition~\ref{prop:influence} gives a sum $\ge16\beta/\pi\ge1$. \qed

\section{Relation to the literature}\label{sec:lit}

\paragraph{Bhattacharjee--Bhattacharya \cite{BB26} (added in v1.4).} After v1.3 of this note (28 September 2026,
not publicly released) was completed, we learned of the public preprint \cite{BB26}, dated 18 October 2026. Its
Theorem~E (p.~2) reads (text of the PDF): ``Let $G=SU(2)$ and $|\beta|<1/(6(d-1))$. The Wilson theory on $\mathbb Z^d$
has exactly one infinite-volume Gibbs measure, and its truncated correlations of local Lipschitz functions decay
exponentially with every rate $\lambda<-\log 6(d-1)|\beta|$. At $d=4$ the range is $\beta_W<2/9$. No single-link
Wasserstein--Dobrushin argument, in the chordal or the geodesic metric, reaches beyond this range.'' with
$\beta_W=N^2\beta$ (their footnote~4, the normalisation of \cite{SZZ23}); its Theorem~F adds a spectral gap of the
transfer matrix in this range, and its Theorem~7 (p.~9) states that for every $c\in\R^4$ and unit $h\in\R^4$,
``$L(c,h)\le\frac14$, with equality at $c=0$. Equivalently, $W_1(\nu_b,\nu_{b'})\le\frac14|b-b'|$ for all
$b,b'\in\R^4$, and the constant $\frac14$ is the best possible.'' \st{read}. To our knowledge \cite{BB26} is the
first public statement of the range of our Theorem~\ref{thm:main}, with stronger conclusions, and of the
maximum-at-zero-field property of Levin's Remark~3.7 for $N=4$ and \emph{all} fields. We make \emph{no priority
claim}, and we have not checked the proofs of \cite{BB26} in full. What this note provides independently is: the
exact one-link coefficient $\Dch(\kappa)=\E_\kappa[x_\perp^2]$ for $\kappa\le4/3$, with the explicit remainder of
Theorem~\ref{thm:chordal}; the exact geodesic coefficient $8\beta/(9\pi)$ for $\kappa\le3\pi/8$
(Theorem~\ref{thm:onelink}) and the range $\pi/16$; two independent interval-arithmetic certificates; the
Lean-checked final steps of \S\ref{sec:chcap} and \S\ref{sec:cap}; and the transfers of the published constants
of Levin, Royer, Gross, Simon and Shen--Zhu--Zhu to the Wilson link (Table~\ref{tab:comp}). A later version of
the same preprint treats a different object (limits with free boundary conditions, gauge-invariant observables),
which we do not compare with Theorem~\ref{thm:main}. Nissim \cite{Nis25} treats $U(N)$ at sufficiently small
't~Hooft coupling; no numerical $SU(2)$ threshold is stated there.

\paragraph{Shen--Zhu--Zhu \cite{SZZ23}.} Their action is (verbatim LaTeX from the arXiv source, eq.~(1.2),
p.~4 of the arXiv PDF)
\begin{center}\small\verb|\cS (Q) \eqdef N\beta  \Re \sum_{p\in \CP^+_{\Lambda_{L}}} \Tr(Q_p),|\end{center}
i.e.\ $\mathcal S(Q)=N\beta\operatorname{Re}\sum_{p\in\mathcal P^+}\operatorname{Tr}(Q_p)$, one term per plaquette; their
Assumption~1.1 is equivalent to (1.3), p.~4: $|\beta|<\frac{1}{16(d-1)}$ for $G=SU(N)$. For $SU(2)$,
$2\beta_{\mathrm{tH}}\ReTr=4\beta_{\mathrm{tH}}\cdot\frac12\ReTr$, so $\beta_W=4\beta_{\mathrm{tH}}$ and their range is
$\beta_W<1/12\approx0.0833$ in $d=4$. Under it they prove (Theorem~1.2, p.~5) uniqueness of the invariant
measure of the Langevin dynamics and convergence of all tight limits of the finite-volume measures, together
with a log-Sobolev inequality (Theorem~1.4) and exponential decay of correlations (Corollary~1.6).
On p.~6 (arXiv v1, the only arXiv version; the published text was not compared) they write (verbatim): ``We remark that uniqueness for small $\beta$ could possibly also be proven
using the method of Dobrushin, see e.g.\ [Dob70]. To this end one would also need to consider the related
Wasserstein metric with respect to the Riemannian distance similarly as we do in this paper. However as we
understand such an argument has not been carried out in detail for lattice Yang Mills in the literature.''
The present note carries out this programme for the gauge model, for uniqueness only (with the Riemannian
distance in Theorem~\ref{thm:geo}, and with the Euclidean distance of $\R^4$ in Theorem~\ref{thm:main}). The
single-site ingredient ($W_1$ on a sphere, derivative along a path of fields written as a covariance) is not
new: it is the method of Levin \cite{Lev80} (Euclidean metric) and Royer \cite{Roy86} (geodesic metric) for
rotator models, and Levin's many-body version already covers the Wilson model (next paragraphs). Neither paper
mentions gauge theories.
Our uniqueness statement is directly of DLR type; we do \emph{not} claim that \cite{SZZ23} lacks DLR
uniqueness (their statement is formulated for the Langevin invariant measure and limits with periodic,
Dirichlet ``or other'' boundary conditions).

\paragraph{Levin: the Euclidean $W_1$ criterion on spheres, many-body interactions.} Levin \cite{Lev80}
(read on the page images of the Project Euclid scan) considers ``the lattice model on $\Z^\nu$ with single
spin space $S^{N-1}$, $N\geqslant2$'' with ``a priori measure $\mu_0$ the invariant probability measure on
$S^{N-1}$'' (p.~65) and the ``Vaserstein distance'' $V(\mu,\nu)=\sup_{\|f\|_L=1}|\mu(f)-\nu(f)|$ over
functions of Lipschitz norm $1$ for ``$\|\cdot\|_2$ [\ldots] Euclidean norm on $\R^N$'' ((3.3)--(3.4), p.~67),
i.e.\ the \emph{chordal} metric of $S^{N-1}\subset\R^N$. His Theorem~3 (p.~67) reads: ``For $\mathbf h,\mathbf h'\in\R^N$,
let $\mu_{\mathbf h},\mu_{\mathbf h'}$ be probability measures on $S^{N-1}$ given by
$d\mu_{\mathbf h}(\mathbf x)=\frac1{Z_{\mathbf h}}e^{\mathbf h\cdot\mathbf x}d\mu_0(\mathbf x)$ [\ldots] Then,
$V(\mu_{\mathbf h},\mu_{\mathbf h'})\leqslant\frac{\sqrt5}{N}\|\mathbf h'-\mathbf h\|_2$. (3.8)'', proved by writing the
distance ``as the variation of an expectation along an appropriate path in $\mathbf h$-space'' (p.~68), with two
lemmas controlling ``separately variation parallel to the `magnetic field', and perpendicular to the `magnetic
field','' the perpendicular one being Lemma~3.2 (p.~68, quoted in \S\ref{sec:chrefl}). Remark~3.6 (p.~73) extends
the uniqueness theorem to ``many-body interactions of the
form $\sum J_X^{(i_a)}\prod_{a\in X}\sigma_a^{i_a}$'' with the condition
``$T>T_C^N\equiv\big(\sup_{a}\sum_{X\ni a}(|X|-1)\mathscr J_X\big)\frac{\sqrt5}{N}$, (3.49)'', where $\mathscr J_X$ is the
supremum of the interaction over configurations; the details are ``[left] to the reader''. Remark~3.7 (p.~73)
conjectures that ``the numerical factor of $\sqrt5$ appearing in Theorem~1 can be replaced by $1$'' and notes that
``the conjecture is true if one shows $\sup_{\|f\|_L=1}|\mu_{\mathbf g}[(\mathbf x\cdot\hat{\mathbf c}-\langle\mathbf x\cdot\hat{\mathbf c}\rangle_{\mathbf g})f]|$ [\ldots]
attains its maximum value at $\mathbf g=0$''. \st{read}. The paper does not mention gauge theories.

\emph{Transfer to the Wilson link} (our computation from the published statement, not a published
result). As noted in \S\ref{sec:influence}, the Wilson weight $e^{(\beta/2)\ReTr U_p}=e^{\beta\langle u_1u_2\bar u_3\bar u_4,\mathbf 1\rangle}$
is of Levin's form (3.48) with $N=4$, $|X|=4$ and $\mathscr J_X=\beta$; each link lies in $6$ plaquettes, so (3.49)
at unit temperature reads $6\cdot3\cdot\beta\cdot\sqrt5/4<1$, i.e.\ $\beta_W<2\sqrt5/45\approx0.09938$.
\st{proved} from \st{read} (exact arithmetic; log \path{calc/levin_transfer.log}).
This already exceeds $1/12$.

\emph{What Theorem~\ref{thm:chordal} adds.} The supremum in Levin's Remark~3.7 is our $\Dch(|\mathbf g|)$ for $N=4$.
Theorem~\ref{thm:chordal} shows that it attains its maximum at $\mathbf g=0$ among $|\mathbf g|\le4/3$, with value
$1/N=1/4$, and identifies it exactly ($\E_{\mathbf g}[x_\perp^2]$, attained by the coordinate orthogonal to the
field). It therefore confirms, for $N=4$ and $|\mathbf g|\le4/3$, the maximum-at-$\mathbf g=0$ property conjectured
in \cite[Rem.~3.7]{Lev80}; it does not prove the conjecture, which is stated for all $N$, all $\mathbf g$ and general
interactions. Since the fields that occur for $\beta<2/9$ satisfy $|\mathbf g|\le6\beta<4/3$, the Wilson threshold is
Levin's transferred bound with $\sqrt5$ replaced by $1$: $2\sqrt5/45\cdot\sqrt5=2/9$ exactly.

\paragraph{Royer: the geodesic $W_1$ criterion on spheres.} Royer \cite{Roy86} (read on the Numdam page
images) works with ``the unit sphere $S^n$ with its natural (geodesic) distance $r$'' and Dobrushin's
distance ``$\mathrm R(\nu_1,\nu_2)=\sup\{\int f\,d\nu_1-\int f\,d\nu_2;\ \mathscr L(f)\leqslant1\}$'' (p.~30).
He defines $C(n,h)$ as ``the best constant in the inequality $\operatorname{var}(f,\nu(h))\leqslant C(n,h)(\mathscr L(f))^2$''
and $C(n)=\sup_hC(n,h)$ (Lemma~1.2, p.~31), proves (3)--(4) quoted in \S\ref{sec:influence}, and bounds $C(n)$:
the proof of Corollary~2.3 (p.~34) gives ``$C(n,h)\leqslant\chi(n,h)<(n-2)^{-1}$; thus
$C(n)\leqslant(n-2)^{-1}$'', Corollary~3.5 (p.~37) ``$C(n,0)=v_n$'', and Theorem~3.6 (p.~37) ``For $n>1$,
$C(n)<v_n+v_{n-1}$'', with $v_n+v_{n-1}=2\sum_{k\geqslant n}k^{-2}$ (p.~37). He states that his estimates
``improve on that given by Levin [10], namely $N/\sqrt5$, for $N\geqslant3$'' (p.~30), and on p.~38 he conjectures
the geodesic analogue of Levin's Remark~3.7: ``one might think that $C(n,h)\leqslant C(n,0)$''. \st{read}.

\emph{Transfer to the Wilson link} (our computation from the published statements, not a published
result). By (3) with $n=3$, $K=\beta$ and $z$ the staple moved by the neighbour, each neighbour has influence
at most $\beta(C(3)/4)^{1/2}$, hence uniqueness when $18\beta(C(3)/4)^{1/2}<1$. With $C(3)\le1$ (Cor.~2.3) this
gives $\beta_W<1/9\approx0.1111$; with $C(3)<v_3+v_2=\pi^2/3-5/2$ (Thm~3.6) it gives
$\beta_W<\big(9\sqrt{\pi^2/3-5/2}\big)^{-1}\approx0.12502$. \st{proved} (exact arithmetic; SymPy logs
\path{calc/royer_transfer.log} and \path{calc/chiffres_v1.log}). At zero field, Royer's route gives $\beta(v_3/4)^{1/2}$ with
$(v_3/4)^{1/2}=\tfrac12(\pi^2/12-\tfrac12)^{1/2}\approx0.28393$, only $0.35\%$ above our exact
$D(0)=8/(9\pi)\approx0.28294$; but this is not uniform in the field and gives no uniqueness statement.

\emph{Where the gain comes from.} Our main threshold $2/9$ exceeds the transferred Royer bound $0.12502$ by a
factor $1.777$, Levin's transferred bound by exactly $\sqrt5\approx2.236$, and $1/12$ by $8/3$; the geodesic
threshold $\pi/16$ exceeds $0.12502$ by $1.5705$, and $2/9$ exceeds $\pi/16$ by $13.2\%$. The gain is not in the
constants at $h=0$: it comes from computing the exact single-link coefficient as a function of $\kappa=\beta|s|$ and
showing that it is maximal at $\kappa=0$ on $[0,6\beta]$ (Theorems~\ref{thm:chordal} and~\ref{thm:onelink}), which
replaces the field-uniform constants $\sqrt5/N$ (Levin) and $C(n)$ (Royer). This maximality is the property
conjectured by Levin (Euclidean metric, all $N$, all fields) and, for the variance constant, by Royer; we establish
it only for $N=4$ and the fields that occur, for the Euclidean $W_1$ coefficient ($\kappa\le4/3$) and for the
geodesic one ($\kappa\le3\pi/8$), by computer-assisted proofs. For $N=4$ the Euclidean property is proved for all
fields in \cite[Theorem~7]{BB26}; we did not check whether \cite{BB26} also covers the geodesic coefficient of
Theorem~\ref{thm:onelink}, and we found no other published proof of either property (literature search not
exhaustive).

\begin{remark}[Numerical evidence and quasi-optimality of the field]\label{rem:chordal}
Before Theorem~\ref{thm:chordal} was proved, linear programs for $\sup_f$ (reduction to functions of two
coordinates, HiGHS; logs \path{calc/lp2d_m14.log} and \path{calc/check1d.log}) gave for
$\max_\alpha\Dch(\kappa,\alpha)$ the values $0.2500$, $0.2474$, $0.2402$, $0.2332$, $0.2164$ at $\kappa=0$, $0.5$, $1$,
$4/3$, $2$,
the maximum over directions being attained for $v\perp\hat s$; the same code reproduces the geodesic $D(0)$ to
$0.15\%$ \st{numerical}. These values agree with $\E_\kappa[x_2^2]=0.25,\,0.24744,\,0.24019,\,0.23332,\,0.21656$ (the
discretised programs are neither upper nor lower bounds; the value $0.2164<0.21656$ at $\kappa=2$ is a
discretisation error, since $\E_\kappa[x_2^2]$ is an exact lower bound for $\Dch(\kappa)$). The full flux bound of
the explicit field \eqref{eq:field} is within $1.7\cdot10^{-4}$ of the linear programs for all directions tested,
so the field is close to optimal \st{numerical}. Beyond $\kappa=4/3$ nothing is certified; the ratio
$\max(J_0,J_1)/(\E x_2^2-\Var x_1)$ grows from $0.32$ at $\kappa=4/3$ to $0.47$ at $\kappa=2$ \st{numerical}.
\end{remark}

\paragraph{Total variation.} The classical total-variation form of Dobrushin's criterion for many-body
interactions on a compact single-spin space is due to Gross \cite{Gro79} and Simon \cite{Sim79}, neither of
whom mentions gauge theories. Gross \cite{Gro79} (read on the page images) treats ``$X$ a compact metric
space'', with ``the cases of principle current interest [\ldots] $X=$ finite set, $S^m$, compact group, or a
finite product of these'' (p.~10), coefficients $\varrho_{a,b}=\frac12\sup\|\mu_b(\cdot|s)-\mu_b(\cdot|t)\|_{\mathrm{Var}}$
((2.3), p.~11), and proves (Theorem~3, p.~22): ``If $\alpha_0<1$ then $\varphi$ has a unique Gibbs state $\sigma$'' with
the decay bound (4.9), where $\alpha_0=2e^{4\|\varphi\|_1}\sup_a\sum_{b\neq a}\zeta(a,b)e^{d(a,b)}$ (4.8), and
``$\|\mu_a(\cdot|t)-\mu_a(\cdot|s)\|_{\mathrm{Var}}\leqslant e^{4\zeta(a,b)}-1$'' ((4.11), Lemma~4.1). \st{read}.
Simon \cite{Sim79} (read on the page images) proves: ``If $\sum_{0\in X}(|X|-1)\|\Phi(X)\|_\infty<1$, then (2)
holds'' (Theorem, p.~184), with ``$\|\Phi(X)\|_\infty$ [replaceable] by $\frac12[\max(\Phi(X))-\min(\Phi(X))]$''
(Remark~1, p.~185), via the Lemma ``$\|\mu_h-\mu_g\|\leqslant\|h-g\|_\infty$'' whose proof writes
$\frac{d}{d\theta}\nu_\theta(f)$ as a covariance ((3)--(4), p.~184); in Remark~5 (p.~185) he observes that a
one-site quantity ``takes its maximum at $a=0$ by the GHS inequality'' for spins in $[-1,1]$. \st{read}.
\emph{Transfer to the Wilson link} (ours): with $\Phi(p)=-\frac\beta2\ReTr U_p$, Simon's condition with his
Remark~1 reads $18\beta<1$, i.e.\ $\beta_W<1/18\approx0.056$; Gross's Theorem~3 gives $36\beta e^{24\beta}<1$, i.e.\
$\beta_W<0.0180$, and his (4.11) in his Theorem~1 gives $\beta_W<\ln(10/9)/4\approx0.0263$ \st{proved} from
\st{read} (log \path{calc/gross_transfer.log}). Chatterjee \cite{Cha20} (read in
arXiv:1811.09770v3) applies the total-variation criterion to $\mathbb Z_2$ (Ising) lattice gauge theory in
$d=4$: ``The easiest way to prove exponential decay of correlations at strong coupling is by using Dobrushin's
condition'' (p.~21) and, in the proof of Theorem~5.2 (p.~24), ``There are at most $18$ such neighbors'', with a
total-variation bound $L_0\lambda$ per neighbour, $L_0$ a universal constant that is not made explicit. \st{read}.
For the $SU(2)$ single-link conditionals the exact one-link total-variation coefficient is
$c_{\mathrm{TV}}(\beta)=\sup\|q_s-q_{s'}\|_{\mathrm{TV}}=\tfrac{4\beta}{3\pi}+O(\beta^2)$, the supremum being over
$|s|,|s'|\le6$ differing by one staple, so $|s-s'|\le2$. \st{proved}: \emph{upper bound} (this is the route
(3) of Simon's Lemma, without the Schwarz step), along the segment
$s(t)=s+t\delta$, $\delta=s'-s$, $|s(t)|\le6$, the density $h_t$ of $q_{s(t)}$ w.r.t.\ Haar measure satisfies
$e^{-12\beta}\le h_t\le e^{12\beta}$, and $\frac{d}{dt}q_{s(t)}=\beta(\langle u,\delta\rangle-\E_{q_{s(t)}}\langle u,\delta\rangle)q_{s(t)}$,
so $\|q_s-q_{s'}\|_{\mathrm{TV}}\le\frac\beta2\sup_t\big(\E_{q_{s(t)}}|\langle u,\delta\rangle|+|\E_{q_{s(t)}}\langle u,\delta\rangle|\big)
\le\beta\big(e^{12\beta}\tfrac{4}{3\pi}+e^{12\beta}-1\big)$, using $\E_{\mathrm{Haar}}|\langle u,e\rangle|=4/(3\pi)$ and
$\E_{\mathrm{Haar}}\langle u,\delta\rangle=0$; \emph{lower bound}, take $s=0$ and $s'=-2w$ (flip the staple
$w\mapsto-w$, the other five summing to $-w$), then $\|q_0-q_{s'}\|_{\mathrm{TV}}=\frac12\E_{\mathrm{Haar}}|e^{\beta\langle u,s'\rangle}/Z-1|
=\frac\beta2\E_{\mathrm{Haar}}|\langle u,s'\rangle|+O(\beta^2)=\frac{4\beta}{3\pi}+O(\beta^2)$. The resulting sum
$24\beta/\pi$ gives the threshold $\approx\pi/24\approx0.131$ \emph{at leading order only}; the finite-$\beta$
threshold, from a floating-point evaluation of $c_{\mathrm{TV}}(\beta)$, is $\approx0.1310$ \st{numerical}.
The ratio of the first-order coefficients is $\frac{12/(9\pi)}{8/(9\pi)}=\frac32$ for the geodesic metric, and
$\frac{4/(3\pi)}{1/4}=\frac{16}{3\pi}\approx1.70$ for the chordal one; the ``gain of $W_1$ over total variation''
therefore depends on the metric.

\paragraph{Summary: five levels.} (i) Total variation: transferred published constants give $\beta_W<1/18$
(Simon) \st{proved} from \st{read}; with the exact one-link coefficient, $\beta_W\approx\pi/24\approx0.131$
(leading order \st{proved}, value \st{numerical}). (ii) Euclidean $W_1$ with Levin's published constant
transferred to the Wilson link: $\beta_W<0.09938$ \st{proved} from \st{read}. (iii) Geodesic $W_1$ with Royer's
published constants transferred: $\beta_W<0.12502$ (Thm~3.6) or $<1/9$ (Cor.~2.3) \st{proved} from \st{read}.
(iv) Geodesic $W_1$ with the exact coefficient, uniform in $|h|\le6\beta$ (Theorem~\ref{thm:geo}):
$\beta_W<\pi/16\approx0.19635$. (v) Euclidean $W_1$ with the exact coefficient, uniform in $|h|\le6\beta$
(Theorem~\ref{thm:main}): $\beta_W<2/9\approx0.2222$. Levels (ii)--(v) exceed $1/12$ rigorously; the
exact-coefficient TV level does so only at leading order and numerically, since the explicit upper bound on
$c_{\mathrm{TV}}$ written above proves the total-variation criterion only for $\beta<0.0414$.

\paragraph{Bakry--\'Emery cannot pass $\beta_W=1/8$.}
\begin{proposition}\label{prop:piflux}
Equip each factor $SU(2)$ with a bi-invariant metric (any constant scale per link). At the $\pi$-flux
configuration ($U_\ell=\pm\mathbf 1$ with every $U_p=-\mathbf 1$; it exists on $\Z^4$ and on periodic boxes of
even side) the quadratic form $\mathrm{Ric}+\beta\,\mathrm{Hess}\,S_W$, $S_W=\sum_p(1-\frac12\ReTr U_p)$, is not
positive definite for $\beta>1/8$ (and, on even periodic boxes, has the eigenvalue $0$ at $\beta=1/8$, so it
is not positive definite there either). Hence the global Bakry--\'Emery
condition $\mathrm{Ric}+\beta\,\mathrm{Hess}\,S_W\ge K>0$ on link variables fails for $\beta_W\ge1/8$. \st{proved}
\end{proposition}
\begin{proof}
In the Frobenius normalisation ($X_a=i\sigma_a/\sqrt2$), $\exp(sX_a)$ has rotation angle $s/\sqrt2$, so
$SU(2)$ is the round $S^3$ of radius $\sqrt2$ and $\mathrm{Ric}=g$. Both $\mathrm{Ric}$ and $\mathrm{Hess}\,S_W$ are
$(0,2)$-tensors unchanged by constant rescalings of the factors, so the sign question is scale-free.
Perturb $U_\ell\mapsto e^{X_\ell}U_\ell$. At a configuration with all $U_\ell$ central ($\pm\mathbf 1$),
$U_p\mapsto U_p\,e^{X_1}e^{X_2}e^{-X_3}e^{-X_4}$, so the second-order expansion of $1-\frac12\ReTr U_p$ at
$U_p=-\mathbf 1$ is minus the one at $U_p=\mathbf 1$, and both configurations are critical. Therefore
$\mathrm{Hess}\,S_W(\pi\text{-flux})=-\mathrm{Hess}\,S_W(\text{flat})=-\frac12\,d^*d\otimes I_3$, where $d^*d=B^{\mathsf T}B$
with $B$ the lattice curl (since $1-\frac12\ReTr e^{Y}=|Y|^2/4+O(|Y|^3)$ and the commutator terms of $Y$ are
traceless). By Fourier analysis the largest eigenvalue of $d^*d$ is $4d=16$, attained at momentum
$(\pi,\pi,\pi,\pi)$ with $d-1=3$ transverse polarisations. So $\lambda_{\min}(\mathrm{Hess}\,S_W)=-8$ with
multiplicity $9$ on even periodic boxes (on $\Z^4$: infimum of the quadratic form, via Weyl sequences), and
$\mathrm{Ric}+\beta\,\mathrm{Hess}\,S_W$ has the eigenvalue $1-8\beta$.
\end{proof}
Checks \st{numerical}: an independent arbiter script (finite differences, $L=2$) finds
$\lambda_{\min}=-8.000000$, multiplicity $9$, $\max|H(\pi)+H(\text{flat})|=7.5\cdot10^{-8}$; a second script
written for this note finds $\max\operatorname{eig}(B^{\mathsf T}B)=16$ with multiplicity $3$ on $L=2$ and $L=4$, and
$Q_\pi(v)=-Q_{\text{flat}}(v)=-\frac12v^{\mathsf T}(B^{\mathsf T}B\otimes I_3)v$ to $2\cdot10^{-7}$ along random directions
on $L=4$. Whether $\sup_U(-\lambda_{\min}\mathrm{Hess}\,S_W)=8$ (so that global Bakry--\'Emery holds up to exactly
$1/8$) is \st{conjecture}. Consistently, \cite{SZZ23} bound $|\mathrm{Hess}|$ by $12\beta_W$ and obtain $1/12$.
Since $2/9>\pi/16>0.12502>1/8$ (the transferred Royer bound exceeds $1/8$ by a margin of $2\cdot10^{-5}$; and,
at leading order only, $\pi/24>1/8$), Dobrushin-type criteria, which need no convexity, go beyond the ceiling of
the global Bakry--\'Emery route on link variables with bi-invariant product metrics. This does not contradict
\cite{CNS25}, which obtains an area law up to $\beta_W<1/6$ by applying Bakry--\'Emery to a different (vertex)
spin system.

\paragraph{Other works.} Osterwalder--Seiler \cite{OS78} prove, by cluster expansion, decay, a mass gap and
uniqueness of the infinite-volume limit for ``$g_0$ sufficiently large'', without an explicit constant
\st{read}. Chatterjee \cite{Cha21} (arXiv:2006.16229v5, p.~7, after Definition~2.3; the published version was not
compared) states that exponential decay under arbitrary boundary conditions ``can be shown to hold for
essentially any lattice gauge theory when $\beta$ is small enough, using the methods of Dobrushin and Shlosman''
\cite{DS85}, without a constant \st{read}. Cao--Nissim--Sheffield
\cite{CNS25} (arXiv, not yet refereed) prove an area law on the torus for $\beta_{\mathrm{tH}}<\frac1{8(d-1)}$,
i.e.\ $\beta_W<1/6$; this is a different statement and is not compared with ours.
Seiler's monograph \cite{Sei82} could not be obtained and was \emph{not consulted}; for Georgii's book
\cite{Geo88} only the table of contents was read (Chapter~8 ``Uniqueness'', \S8.1 ``Dobrushin's condition of
weak dependence'', p.~140), and its content was \emph{not consulted}; Simon's book \cite{Sim93} was \emph{not
consulted}. No claim of this note rests on any of them (see \S\ref{sec:limits}, item~4).

\begin{table}[htbp]
\centering\small
\begin{tabularx}{\linewidth}{@{}>{\raggedright\arraybackslash}p{0.25\linewidth}>{\raggedright\arraybackslash}X>{\raggedright\arraybackslash}p{0.19\linewidth}>{\raggedright\arraybackslash}p{0.21\linewidth}@{}}
\toprule
source & statement & range in $\beta_W$ ($SU(2)$, $d=4$) & status\\
\midrule
\cite{OS78} & decay, gap, unique limit & small, not explicit & \st{read}\\
\cite{Cha20} & TV-Dobrushin, decay, $\mathbb Z_2$ gauge, $d=4$ & not explicit ($L_0$) & \st{read}; other group\\
\cite{Gro79} Thm~3, transferred (here) & uniqueness, decay & $<0.0180$ & \st{proved} from \st{read}\\
\cite{Sim79} Thm + Rem.~1, transferred (here) & uniqueness & $<1/18\approx0.056$ & \st{proved} from \st{read}\\
\cite{SZZ23} & uniqueness, LSI, decay & $<1/12\approx0.083$ & \st{read}\\
\cite{Lev80} Rem.~3.6, transferred (here) & uniqueness & $<2\sqrt5/45\approx0.0994$ & \st{proved} from \st{read}\\
\cite{Roy86} Cor.~2.3, transferred (here) & uniqueness & $<1/9\approx0.111$ & \st{proved} from \st{read}\\
Prop.~\ref{prop:piflux} & global BE on links fails & $\ge1/8=0.125$ & \st{proved}\\
\cite{Roy86} Thm~3.6, transferred (here) & uniqueness & $<0.12502$ & \st{proved} from \st{read}\\
TV-Dobrushin, exact one-link coefficient (here) & uniqueness & $\approx\pi/24\approx0.131$ & leading order \st{proved}, value \st{numerical}\\
\cite{CNS25} & area law on torus (different) & $<1/6\approx0.167$ & \st{read}, preprint\\
Theorem~\ref{thm:geo} (geodesic $W_1$, exact coefficient) & uniqueness (DLR) only & $<\pi/16\approx0.196$ & see Thm~\ref{thm:geo}\\
Theorem~\ref{thm:main} (Euclidean $W_1$, exact coefficient) & uniqueness (DLR) only & $<2/9\approx0.222$ & see Thm~\ref{thm:main}\\
\cite{BB26} Thms~E, F (preprint, 18 Oct.\ 2026) & uniqueness, decay, transfer-matrix gap & $<2/9\approx0.222$ & \st{read}; proofs not checked here\\
\bottomrule
\end{tabularx}
\caption{Explicit ranges, all converted to $\beta_W$ for $SU(2)$ in $d=4$ (Table~\ref{tab:conv}). Rows marked
``transferred (here)'' apply published statements of \cite{Gro79,Sim79,Lev80,Roy86} to the Wilson link;
they are not published results, and none of these papers mentions gauge theories.}
\label{tab:comp}
\end{table}

\section{Limitations}\label{sec:limits}
\begin{enumerate}
\item \emph{Uniqueness only.} No decay of correlations, spectral gap, log-Sobolev inequality or area law is
claimed at $\beta<2/9$ in this note. The preprint \cite{BB26} states exponential decay and a transfer-matrix gap
in the same range (proofs not checked here); among refereed publications, the best explicit range we know for
these conclusions remains $1/12$ \cite{SZZ23}. Explicit decay statements covering four-body plaquette interactions do exist in total
variation: Gross's Theorem~3 \cite{Gro79}, transferred, gives uniqueness and exponential decay for
$\beta_W<0.0180$, far below $1/12$. Gross's Theorem~1 \cite[p.~11]{Gro79} turns any total-variation Dobrushin sum
below $1$ into exponential decay (take $d(a,b)=r|a-b|$ with $r>0$ small); with a certified bound on $c_{\mathrm{TV}}$
it would give decay up to the TV threshold $\approx0.131$, hence beyond $1/12$ \st{conjecture} (not certified,
not written out). For the $W_1$ criterion, Levin \cite[Remark~3.5, p.~72]{Lev80} states that, ``parallel to the
discussion of Gross'', the truncated spin--spin correlations decay exponentially under the $W_1$ condition with a
semi-metric weight (his (3.46)--(3.47)); applied to our influence matrices this would plausibly give exponential
decay of correlations of Lipschitz observables for $\beta<2/9$ \st{conjecture}. Nothing is claimed.
\item \emph{Formalisation boundary.} For the chordal certificate, the flux bound, the explicit field and the
elimination of the direction (\S\S\ref{sec:chfield}--\ref{sec:chalpha}) are proved on paper; the box bounds, the
positivity of $m$ and the constant $0.0084$ are interval-certified; Lean checks only the final rational
implication, two algebraic steps and the threshold arithmetic. For the geodesic certificate, the chain
``integrals $\to$ coefficient bounds'' (Arb, reduction of \S\ref{sec:reduction}) is outside Lean, and so are the
constant $0.0085$ and the covering of $[0,\pi]$ by the boxes (\S\ref{sec:cap}). Kantorovich--Rubinstein duality is
quoted from \cite[Thm~10.8, p.~53]{DS25}; McShane's extension, the divergence theorem for Lipschitz fields, the
Holley--Stroock perturbation argument \cite{HS87} and the weak solvability of $-L\varphi=R$ (Lax--Milgram +
Poincar\'e) are standard and written out here; none of them is formalised.
\item \emph{History and audit of the certificates.} The first geodesic certificate was wrong (see the erratum);
the present one passed guards designed to catch that failure mode and two independent recomputations. The chordal
certificate was written with the same guards and is confirmed by an independent implementation with different
arithmetic, tail bounds and quadrature (\S\ref{sec:chcap}); its mathematical reduction was re-derived line by line
in that independent verification. Neither certificate, and none of the proofs, has been audited or refereed by a
party outside the project.
\item \emph{Novelty and prior art.} The single-site $W_1$ machinery on spheres is due to Levin \cite{Lev80}
(Euclidean metric, many-body multilinear interactions, which cover the Wilson model; transferred, $0.0994$)
and Royer \cite{Roy86} (geodesic metric; transferred, $0.12502$). The field-uniformity of the one-site
coefficient (maximum at zero field) was conjectured by both, and the total-variation criterion for many-body
interactions is due to Gross \cite{Gro79} and Simon \cite{Sim79}. None of these papers mentions gauge theories,
and none contains $\beta<\pi/16$, $\beta<2/9$ or a better threshold. The contribution of this note is the exact
field-uniform single-link coefficient for the Euclidean and for the geodesic metric, its certification, and the
resulting thresholds $2/9$ and $\pi/16$.
\emph{Limitation of the literature search:} Seiler's monograph \cite{Sei82} (not obtainable), the content of
Georgii's book \cite{Geo88} (table of contents only) and Simon's book \cite{Sim93} were not consulted; any of them
may contain a Dobrushin-type statement for plaquette interactions or for $SU(2)$ gauge fields. No claim of this
note rests on them, but a priority claim for the gauge application should be read with this limitation in mind.
\item \emph{Range of fields.} The chordal identity is certified only for $\kappa\le4/3$, which is exactly what
$\beta<2/9$ requires; whether $\Dch(\kappa)\le1/4$ for all $\kappa$ (Levin's conjecture for $N=4$) is not
addressed. Likewise the geodesic bound is certified only for $\kappa\le59/50$.
\item \emph{The value $\pi/24$} is exact only at leading order in $\beta$; the finite-$\beta$ TV threshold is a
floating-point estimate, and has not been recertified in interval arithmetic.
\item \emph{Physics.} The result is a strong-coupling statement ($g^2>18$, with $\beta_W=4/g^2$). It says nothing
about the continuum limit or the weak-coupling regime relevant to the Yang--Mills mass gap problem.
\end{enumerate}

\begin{thebibliography}{99}
\bibitem{Dob70} R.~L.~Dobrushin, \ru{Задание системы случайных величин при помощи условных распределений},
Teor.\ Veroyatn.\ Primen.\ 15:3 (1970) 469--497; English transl.\ \emph{Prescribing a system of random variables
by conditional distributions}, Theory Probab.\ Appl.\ 15 (1970) 458--486, doi:10.1137/1115049.
\bibitem{SZZ23} H.~Shen, R.~Zhu, X.~Zhu, \emph{A stochastic analysis approach to lattice Yang--Mills at strong
coupling}, Comm.\ Math.\ Phys.\ 400 (2023) 805--851, doi:10.1007/s00220-022-04609-1, arXiv:2204.12737.
(Page numbers above refer to the arXiv v1 PDF, the only arXiv version.)
\bibitem{OS78} K.~Osterwalder, E.~Seiler, \emph{Gauge field theories on a lattice}, Ann.\ Phys.\ 110 (1978)
440--471, doi:10.1016/0003-4916(78)90039-8.
\bibitem{Sei82} E.~Seiler, \emph{Gauge Theories as a Problem of Constructive Quantum Field Theory and
Statistical Mechanics}, Lecture Notes in Physics 159, Springer (1982), doi:10.1007/3-540-11559-5.
\emph{Not consulted.}
\bibitem{CNS25} S.~Cao, R.~Nissim, S.~Sheffield, \emph{Dynamical approach to area law for lattice Yang--Mills},
arXiv:2509.04688v2 (2025).
\bibitem{Cha21} S.~Chatterjee, \emph{A probabilistic mechanism for quark confinement}, Comm.\ Math.\ Phys.\ 385
(2021) 1007--1039, doi:10.1007/s00220-021-04086-y, arXiv:2006.16229. (Page number above refers to the arXiv v5
PDF; the published version was not compared.)
\bibitem{DS85} R.~L.~Dobrushin, S.~B.~Shlosman, \emph{Constructive criterion for the uniqueness of Gibbs field},
in: Statistical Physics and Dynamical Systems, Birkh\"auser (1985) 347--370, doi:10.1007/978-1-4899-6653-7\_20.
\bibitem{DS25} T.~C.~Dorlas, B.~Savoie, \emph{On the Wasserstein distance and the Dobrushin uniqueness theorem},
Rev.\ Math.\ Phys.\ 37 (2025), doi:10.1142/S0129055X25300018, arXiv:2305.19371. (Theorem and page
numbers refer to the arXiv v4 PDF, 66 pp.; the published version was not compared.)
\bibitem{HS87} R.~Holley, D.~Stroock, \emph{Logarithmic Sobolev inequalities and stochastic Ising models},
J.\ Stat.\ Phys.\ 46 (1987) 1159--1194, doi:10.1007/BF01011161.
\bibitem{Roy86} G.~Royer, \emph{Study of Dobrushin's critical coupling in rotator models}, Ann.\ Inst.\
H.~Poincar\'e Phys.\ Th\'eor.\ 44:1 (1986) 29--38 (Numdam \path{AIHPA_1986__44_1_29_0}; metadata from Numdam,
no Crossref DOI; read on the page images, pp.~30, 31, 34, 37).
\bibitem{Lev80} S.~L.~Levin, \emph{Application of Dobrushin's uniqueness theorem to $N$-vector models},
Comm.\ Math.\ Phys.\ 78 (1980) 65--74, doi:10.1007/BF01941969 (Project Euclid \path{cmp/1103908501}; read on the
page images, pp.~65, 66, 67, 68, 71, 72, 73).
\bibitem{Cha20} S.~Chatterjee, \emph{Wilson loops in Ising lattice gauge theory}, Comm.\ Math.\ Phys.\ 377
(2020) 307--340, doi:10.1007/s00220-020-03738-9, arXiv:1811.09770. (Page numbers above refer to the arXiv v3
PDF; the published version was not compared.)
\bibitem{Gro79} L.~Gross, \emph{Decay of correlations in classical lattice models at high temperature},
Comm.\ Math.\ Phys.\ 68 (1979) 9--27, doi:10.1007/BF01562538 (Project Euclid \path{cmp/1103905263}; read on the
page images, pp.~10, 11, 21, 22).
\bibitem{Sim79} B.~Simon, \emph{A remark on Dobrushin's uniqueness theorem}, Comm.\ Math.\ Phys.\ 68 (1979)
183--185, doi:10.1007/BF01418127 (Project Euclid \path{cmp/1103905311}; read on the page images, pp.~183--185).
\bibitem{Sim93} B.~Simon, \emph{The Statistical Mechanics of Lattice Gases, Volume~I}, Princeton University
Press (1993), doi:10.1515/9781400863433 (metadata from Crossref). \emph{Not consulted.}
\bibitem{Geo88} H.-O.~Georgii, \emph{Gibbs Measures and Phase Transitions}, de Gruyter, Berlin--New York (1988),
doi:10.1515/9783110850147. \emph{Content not consulted}; table of contents read (Deutsche Nationalbibliothek
record 880793589).
\bibitem{BB26} D.~Bhattacharjee, U.~Bhattacharya, \emph{Mass Gap of Lattice Yang--Mills Theory up to the Dobrushin
Ceiling}, preprint hexagon:2610.00341v1, dated 18 October 2026, \url{https://hexagonmath.org/pdf/2610.00341v1}.
\bibitem{Nis25} R.~Nissim, \emph{U(N) lattice Yang--Mills in the 't~Hooft regime}, preprint arXiv:2510.22788 (2025), \url{https://arxiv.org/abs/2510.22788}.
\end{thebibliography}

\end{document}
